% Representations of reductive groups over finite fields
% By P. Deligne and G. Lusztig
% Pages 1-50 (journal pages 103-152)
\documentclass[11pt]{amsart}
\usepackage[utf8]{inputenc}
\usepackage[T1]{fontenc}
\usepackage{amsmath,amssymb,amsthm}
\usepackage{enumerate}
\usepackage{mathrsfs}
\usepackage{graphicx}

% Theorem environments
\theoremstyle{plain}
\newtheorem{theorem}{Theorem}[section]
\newtheorem{lemma}[theorem]{Lemma}
\newtheorem{proposition}[theorem]{Proposition}
\newtheorem{corollary}[theorem]{Corollary}
\newtheorem{definition}[theorem]{Definition}
\theoremstyle{remark}
\newtheorem{remark}[theorem]{Remark}
\newtheorem{example}[theorem]{Example}

% Math operators
\DeclareMathOperator{\Hom}{Hom}
\DeclareMathOperator{\Ker}{Ker}
\DeclareMathOperator{\Ind}{Ind}
\DeclareMathOperator{\Tr}{Tr}
\DeclareMathOperator{\tr}{tr}
\DeclareMathOperator{\pr}{pr}
\DeclareMathOperator{\ad}{ad}
\DeclareMathOperator{\rank}{rank}
\DeclareMathOperator{\id}{Id}
\DeclareMathOperator{\End}{End}
\DeclareMathOperator{\Coker}{Coker}
\DeclareMathOperator{\Gr}{Gr}
\DeclareMathOperator{\Tor}{Tor}
\DeclareMathOperator{\Aut}{Aut}
\DeclareMathOperator{\red}{red}
\DeclareMathOperator{\Spec}{Spec}
\DeclareMathOperator{\St}{St}
\DeclareMathOperator{\Res}{Res}
\DeclareMathOperator{\Lie}{Lie}

% Common macros
\newcommand{\GF}{G^F}
\newcommand{\TF}{T^F}
\newcommand{\ZZ}{\mathbb{Z}}
\newcommand{\QQ}{\mathbb{Q}}
\newcommand{\Fq}{\mathbb{F}_q}
\newcommand{\Gal}{\overline{\mathbb{Q}}_\ell}
\newcommand{\Hc}{H_c^*}
\newcommand{\Hci}{H_c^i}
\newcommand{\Hcf}{H_c^i(X, \overline{\mathbb{Q}}_\ell)}
\newcommand{\Grr}{\mathscr{R}}

\begin{document}

\title{Representations of reductive groups over finite fields}
\author{P. Deligne and G. Lusztig}
\begin{abstract}
Let us consider a connected, reductive algebraic group $G$, defined over a finite field $\mathbb{F}_q$, with Frobenius map $F$. We shall be concerned with the representation theory of the finite group $G^F$, over fields of characteristic $0$. In this paper we prove Macdonald's conjecture. More precisely, for $T$ an $F$-stable maximal torus and $\theta$ an arbitrary character of $T^F$ we construct virtual representations $R_T^\theta$ which have all the required properties.
\end{abstract}
\maketitle

\section*{Introduction}

Let us consider a connected, reductive algebraic group $G$, defined over a finite field $\mathbb{F}_q$, with Frobenius map $F$. We shall be concerned with the representation theory of the finite group $G^F$, over fields of characteristic $0$.

In 1968, Macdonald conjectured, on the basis of the character tables known at the time ($GL_n$, $Sp_4$), that there should be a well defined correspondence which, to any $F$-stable maximal torus $T$ of $G$ and a character $\theta$ of $T^F$ in general position, associates an irreducible representation of $G^F$; moreover, if $T$ modulo the centre of $G$ is anisotropic over $\mathbb{F}_q$, the corresponding representation of $G^F$ should be cuspidal (see \textit{Seminar on algebraic groups and related finite groups}, by A. Borel et al., Lecture Notes in Mathematics, 131, pp.\ 117 and 101). In this paper we prove Macdonald's conjecture. More precisely, for $T$ as above and $\theta$ an arbitrary character of $T^F$ we construct virtual representations $R_T^\theta$ which have all the required properties.

These are defined in Chapter 1 as the alternating sum of the cohomology with compact support of the variety of Borel subgroups of $G$ which are in a fixed relative position with their $F$-transform, with coefficients in certain $G^F$-equivariant locally constant $\ell$-adic sheaves of rank one. This generalizes a construction made by Drinfeld for $SL_2$ (see Ch.\ 2).

In Chapter 4 we prove a character formula for $R_T^\theta$, based on the fixed point formula of Chapter 3. This character formula is in terms of certain ``Green functions'' on the unipotent elements; in Chapters 6, 7, 8 we prove that these Green functions satisfy all the axioms predicted by Springer, Kilmoyer and Macdonald (\cite{SK}, \cite{Mac}).

By 6.8, $\varepsilon R_T^\theta$ is irreducible if $\theta$ is in general position and the vanishing theorem (9.9) gives an explicit model for it provided that $q$ is not too small (if $G$ is a classical group or $G_2$, any $q$ will do; in the general case $q > 30$ is sufficient).

In Chapter 10 we study the irreducible components of the Gelfand-Graev representation of $G^F$, assuming that the centre of $G$ is connected. The proof uses the results of Chapter 5 together with the disjointness theorem (6.2).

Finally, in Chapter 11 we discuss the case of the Suzuki and Ree groups.

It would be very desirable to find formulas for the Green functions more explicit than (4.1.2). Such formulas are known for $GL_n$ (Green \cite{Green}), $Sp_4$ (Srinivasan \cite{Sri}) and $G_2$ (Chang, Ree \cite{ChRee}). Very recently, Kazhdan has proved, using results of Springer, that the Green functions can be expressed as exponential sums on the Lie algebra (see \cite{Mac}) provided that the characteristic is good and $q$ is not too small.

Some of the results in this paper were announced in \cite{DL1}, \cite{DL2}.

The second author would like to thank the I.H.E.S.\ for its hospitality during part of the time of preparation of this paper.

\section*{Table of Contents}

\begin{tabular}{ll}
Conventions and standard notations & \dotfill 104 \\
Chapter 1.\ Some basic definitions & \dotfill 105 \\
Chapter 2.\ Examples in the classical groups & \dotfill 116 \\
Chapter 3.\ A fixed point formula & \dotfill 118 \\
Chapter 4.\ The character formula & \dotfill 123 \\
Chapter 5.\ Characters of tori & \dotfill 126 \\
Chapter 6.\ Intertwining numbers & \dotfill 135 \\
Chapter 7.\ Computations on semisimple elements & \dotfill 140 \\
Chapter 8.\ Induced and cuspidal representations & \dotfill 146 \\
Chapter 9.\ A vanishing theorem & \dotfill 147 \\
Chapter 10.\ A decomposition of the Gelfand-Graev representation & \dotfill 154 \\
Chapter 11.\ Suzuki and Ree groups & \dotfill 160
\end{tabular}

\section*{Conventions and standard notations}

\subsection*{0.1} In this paper, $k$ is always an algebraically closed field and $p$ is its characteristic exponent. The following assumptions and definitions will be in force in Chapters 4 to 8, in Chapter 10, and in parts of Chapters 1 and 9.

\textbf{(0.1.1)} $p$ is a prime, and $k$ is an algebraic closure of the prime field $\mathbb{F}_p$ of characteristic $p$. For $q$ a power of $p$, $\mathbb{F}_q$ is the subfield with $q$ elements of $k$. $G$ is a reductive algebraic group over $k$, obtained by extension of scalars from $G_0$ over $\mathbb{F}_q$. We denote by $F$ the corresponding Frobenius endomorphism $F: G \to G$ (as well as the Frobenius endomorphism of any $k$-scheme defined over $\mathbb{F}_q$).

\subsection*{0.2} $\ell$ is a prime different from $p$, and $\overline{\mathbb{Q}}_\ell$ is an algebraic closure of the $\ell$-adic field. When there is no ambiguity on $\ell$, we will write simply $H_c^*(X)$ (resp.\ $H^*(X)$) for $H_c^*(X, \overline{\mathbb{Q}}_\ell)$ (resp.\ $H^*(X, \overline{\mathbb{Q}}_\ell)$) ($\ell$-adic cohomology of $X$, a scheme over $k$). The groups $\mathbb{Z}/\ell^n(1)$ are the groups $\mu_{\ell^n}(k)$ of $\ell^n$-roots of unity in $k$; they form a projective system, with transition maps $x \mapsto x^{\ell^{n-m}}$; the projective limit is $\mathbb{Z}_\ell(1)$. One defines $\mathbb{Z}_\ell(n) = \mathbb{Z}_\ell(1)^{\otimes n}$, and $\mathbb{Z}_\ell(-n) =$ the dual of $\mathbb{Z}_\ell(n)$. The symbol $(n)$ is for a Tate twist: tensoring with $\mathbb{Z}_\ell(n)$. We will make an accidental use of the similarly defined group $\hat{\mathbb{Z}}_{p'}(1) = \varprojlim_{(n,p)=1} \mu_n(k)$.

\subsection*{0.3} For $H$ a finite group, $\mathcal{R}(H)$ is the Grothendieck group of the finite dimensional representations of $H$ over $\overline{\mathbb{Q}}_\ell$. If $f$ and $f'$ are class functions on $H$, with values in a cyclotomic field (often $\subset \overline{\mathbb{Q}}_\ell$), the inner product $\langle f, f' \rangle_H$ (or simply $\langle f, f' \rangle$) is $(1/|H|) \sum_{x \in H} f(x)\overline{f'(x)}$, where $\bar{\ }$ is the automorphism inducing complex conjugation on the roots of unity. The same notation $\langle \ ,\ \rangle$ will be used for the inner product of elements of $\mathcal{R}(H) \otimes \mathbb{Q}$, identified with their characters.

\subsection*{0.4} The length function on a Coxeter group $W$ (with canonical generators $s_1, \dots, s_n$) will be denoted by $l(\ )$. A reduced expression for $w \in W$ is a decomposition $w = s_{i_1} \cdots s_{i_k}$ with $l(w) = k$.

\subsection*{0.5} Miscellaneous:
\begin{itemize}
\item[$\pr_i$:] $i$th projection (as in $\pr_1: X \times X \to X$);
\item[$X^T$:] the fixed point subscheme of the endomorphism $T$ of the scheme $X$ (for instance: $G^F = G_0(\mathbb{F}_q)$);
\item[$p'$ (in index):] away from $p$ (as in $\hat{\mathbb{Z}}_{p'}$, completion away from $p$) or the subgroup of elements of order prime to $p$ (as in $(\mathbb{Q}/\mathbb{Z})_{p'}$);
\item[$\ad g$:] the inner automorphism $x \mapsto gxg^{-1}$, or maps deduced from it;
\item[$1$ or $e$:] the identity element of a group;
\item[$\Tr(f, V^*)$,] for $f$ an endomorphism of a graded vector space, is $\sum (-1)^i \Tr(f, V^i)$;
\item[reductive:] reductive groups are meant to be connected and smooth.
\end{itemize}

The Jordan decomposition of an element $g \in G$ ($G$ as in (0.1.1)) is $g = su$ where $s$ is a semisimple element and $u$ is a unipotent element commuting with $s$.

\subsection*{0.6} We will often identify a scheme over $k$ with the set of its $k$-rational points. This should cause no confusion.

\section{Some basic definitions}

\subsection{1.1} Suppose that in some category we are given a family $(X_i)_{i \in I}$ of objects and a compatible system of isomorphisms $\varphi_{ij}: X_i \xrightarrow{\sim} X_j$. This is as good as giving a single object $X$, the ``common value'' or ``projective limit'' of the family. This projective limit is provided with isomorphisms $\psi_i: X \xrightarrow{\sim} X_i$ such that $\varphi_{ji}\psi_i = \psi_j$. We will use such a construction to define the maximal torus $T$ and the Weyl group $W$ of a connected reductive algebraic group $G$ over $k$.

As index set $I$, we take the set of pairs $(T, B)$ consisting of a maximal torus $T$ and a Borel subgroup $B$ containing $T$. For $i \in I$, $i = (T_i, B_i)$, we take $T_i = T$, $W_i = N(T)/T$. The isomorphism $\varphi_{ij}$ is the isomorphism induced by $\ad g$ where $g$ is any element of $G$ conjugating $i$ into $j$; these elements $g$ form a single right $T_i$-coset, so that $\varphi_{ji}$ is independent of the choice of $g$.

One similarly defines the root system of $T$, its set of simple roots, the action of $W$ on $T$ and the fundamental reflections in $W$.

Let $F: G \to G$ be an isogeny. For any $i \in I$, $i = (T, B)$, $F(i) = (F(T), F(B))$ is again in $I$; $F$ induces an isogeny $F: T_i \to T_{F(i)}$ and an isomorphism $F: W_i \xrightarrow{\sim} W_{F(i)}$. The corresponding endomorphism $\psi_i^{-1}F\psi_i$ of the torus $T$ (resp. automorphism of the Weyl group $W$) is independent of the choice of $i$; we say that it is \textit{induced by} $F$.

\subsection{1.2} Let $X$ (or $X_G$) be the set of all Borel subgroups of $G$. The group $G$ acts on $X$ by conjugation, and $X$ is a smooth projective homogeneous space of $G$. For each Borel subgroup $B$ of $G$, $B$ is the stabilizer of the corresponding point of $X$, hence there is a natural isomorphism $G/B \xrightarrow{\sim} X: g \mapsto gBg^{-1}$.

The set of orbits of $G$ in $X \times X$ can be identified with the Weyl group $W$ of $G$ as follows: for any $i = (T, B)$ as in (1.1), use the composite bijection
\[
W \xleftarrow{\sim} N(T)/T \xleftarrow{\sim} B\backslash G/B \xrightarrow{\sim} G\backslash(X \times X);
\]
this is independent of the choice of $(T, B)$. We denote by $O(w)$ the orbit corresponding to $w \in W$; this is the orbit of $(B, \dot{w}B\dot{w}^{-1})$, where $\dot{w} \in N(T)$ represents $w$. We shall say that two Borel subgroups $B'$, $B''$ of $G$ are in relative position $w$, $w \in W$, if and only if $(B', B'') \in O(w)$. In diagrams, we will picture this by $B' \xrightarrow{w} B''$.

The basic properties of Bruhat decomposition can be expressed as follows:

(a) If $w = w_1 w_2$, with $l(w) = l(w_1) + l(w_2)$ then

(a$_1$) $(B', B'') \in O(w_1)$ and $(B'', B''') \in O(w_2) \Rightarrow (B', B''') \in O(w)$;

(a$_2$) if $(B', B''') \in O(w)$, there is one and only one $B''$ such that $(B', B'') \in O(w_1)$ and $(B'', B''') \in O(w_2)$.

On the scheme level: $O(w_1) \times_X O(w_2) \xrightarrow{\sim} O(w)$.

(b) Let $s$ be an elementary reflection and let $B$ be a Borel subgroup.

(b$_1$) $P_s = B \cup BsB$ is a (minimal parabolic) subgroup, i.e., if $(B', B'') \in O(s)$ and $(B'', B''') \in O(s)$, then either $(B', B''') \in O(s)$ or $B' = B'''$.

(b$_2$) The quotient of $L = P_s/U_{P_s}$ ($U_{P_s} =$ unipotent radical of $P_s$) by its centre is isomorphic to $PGL(2)$. The space $X_L$ is hence a projective line. The inverse image map $X_L \to X$ has as image the set of Borel subgroups $B'$ in relative position $e$ or $s$ with $B$.

(b$_3$) Via either projection, $\overline{O(s)} = O(s) \cup O(e) \subset X \times X$ is hence a fibre space over $X$, with fibre $\mathbb{P}^1$. It is provided with the section $O(e)$; this reduces its structural group from the projective to the affine group. The complement $O(s)$ of that section is a fibre space with fibre the affine line.

\subsection{1.3} The assumptions (0.1.1) are in force in the rest of this chapter. The scheme $X$ is hence defined over $\mathbb{F}_q$ and provided with a Frobenius map $F: X \to X$.

\begin{definition}[1.4]
For $w$ in the Weyl group $W$ of $G$, $X(w) \subset X$ is the locally closed subscheme of $X$ consisting of all Borel subgroups $B$ of $G$ such that $B$ and $F(B)$ are in relative position $w$.
\end{definition}

One can also regard $X(w)$ as the intersection, in $X \times X$, of $O(w)$ with the graph of Frobenius. It is easily checked that this intersection is transverse. The orbit $O(w)$ being smooth of dimension $\dim(X) + l(w)$, it follows that $X(w)$ is smooth and purely of dimension $l(w)$. The subscheme $X(w)$ of $X$ is $G^F$-stable. Hence, for each prime number $\ell \neq p$, $G^F$ acts on the $\ell$-adic cohomology with compact supports of $X(w)$.

\begin{definition}[1.5]
$R^1(w)$ is the virtual representation
\[
\sum_i (-1)^i H_c^i(X(w), \overline{\mathbb{Q}}_\ell)
\]
of $G^F$ (an element of the Grothendieck group of representations of $G^F$ over $\overline{\mathbb{Q}}_\ell$).
\end{definition}

For $w = e$, $X(w)$ is of dimension $0$; it is the set of rational Borel subgroups, and $R^1(w)$ is induced by the unit representation of $B^F$, where $B$ is an $F$-stable Borel subgroup.

It will follow from (3.3) that the character of $R^1(w)$ has integral values, independent of $\ell$. This will justify omitting $\ell$ from the notation; $R^1(w)$ could be also regarded as a complex virtual representation of $G^F$.

An element of the Weyl group $W$ is said to be \textit{$F$-conjugate} to $w \in W$, if it is of the form $w_1 w F(w_1)^{-1}$, for some $w_1 \in W$. We denote by $W^F$ the set of $F$-conjugacy classes in $W$. We shall recall in (1.14), that the $G^F$-conjugacy classes of $F$-stable (i.e., $\mathbb{F}_q$-rational) maximal tori in $G$ are parametrized by $W^F$.

\begin{theorem}[1.6]
$R^1(w)$ depends only on the $F$-conjugacy class of $w$.
\end{theorem}

Let $w$ and $w'$ be $F$-conjugate.

\textit{Case 1.} $w = w_1 w_2$, $w' = w_2 F(w_1)$ and $l(w) = l(w_1) + l(w_2) = l(w_2) + l(F(w_1)) = l(w')$ (0.4). For $B \in X(w)$, there is a unique Borel subgroup $\alpha B$ such that $(B, \alpha B) \in O(w_1)$ and $(\alpha B, FB) \in O(w_2)$; we have the diagram of relative positions:
\[
\begin{array}{ccccc}
B & \xrightarrow{w_1} & \alpha B & \xrightarrow{w_2} & F(B) \\
& & \downarrow{w_2} & & \downarrow{F(w_1)} \\
& & F(\alpha B) & &
\end{array}
\]
(the bottom one because $l(w_2 F(w_1)) = l(w_2) + l(F(w_1))$), and $\alpha B \in X(w')$. By the same argument applied to $w_2 F(w_1)$ and $F(w_1)F(w_2) = F(w)$, we get a map $\beta: X(w') \to X(Fw)$. The diagram
\[
\begin{array}{ccc}
X(w) & \xrightarrow{\alpha} & X(w') \\
\downarrow{F} & & \downarrow{F} \\
X(Fw) & \xrightarrow{\beta} & X(Fw')
\end{array}
\]
is commutative. The vertical maps induce equivalences of \'etale sites, hence so do $\alpha$ and $\beta$ [SGA 1, IX, 4.10]. The resulting isomorphism
\[
H_c^*(X(w')) \xrightarrow{\sim} H_c^*(X(w))
\]
is $G^F$-equivariant, whence (1.6) in this case.

\textit{Case 2.} For some fundamental reflection $s$, we have $w' = s w F(s)$ and $l(w') = l(w) + 2$. For $B \in X(w')$, we can find two Borel subgroups $\gamma B$, $\delta B$ such that $(B, \gamma B) \in O(s)$, $(\gamma B, \delta B) \in O(w)$, $(\delta B, FB) \in O(F(s))$; moreover, $\gamma B$ and $\delta B$ are uniquely determined by these requirements. We define a partition $X(w') = X_1 \cup X_2$ by
\[
X_1 = \{B \in X(w') \mid \delta B = F(\gamma B)\}, \qquad X_2 = \{B \in X(w') \mid \delta B \neq F(\gamma B)\}.
\]
Note that $X_1$ is a closed subscheme of $X(w')$, while $X_2$ is an open one. We have $\gamma: X_1 \to X(w)$ and, for $B' \in X(w)$, $\gamma^{-1}(B')$ is the set of all Borel subgroups $B$ such that $(B, B') \in O(s)$; hence $\gamma^{-1}(B')$ is an affine line over $k$, and $X_1$ is (via $\gamma$) an affine line bundle over $X(w)$. It follows that $\gamma$ induces an isomorphism
\begin{equation}\label{eq:1.6.1}
H_c^i(X_1) \xrightarrow{\sim} H_c^{i-2}(X(w))(-1), \qquad i > 0.
\end{equation}

If $B \in X_2$, we have
\[
(\delta B, FB) \in O(F(s)), \quad (FB, F(\gamma B)) \in O(F(s)), \quad \delta B \neq F(\gamma B)
\]
hence $(\delta B, F(\gamma B)) \in O(F(s))$. Since
\[
(F(\gamma B), F(\delta B)) \in O(F(w)), \quad l(F(sw)) = l(F(s)) + l(F(w)),
\]
it follows that $(\delta B, F(\delta B)) \in O(F(sw))$. Thus we have $\delta: X_2 \to X(F(sw))$. Let $X_2' = \{(B, \bar{B}) \in X_2 \times X(sw) \mid F(\bar{B}) = \delta B\}$ and let $\delta': X_2' \to X(sw)$ be defined by $\delta'(B, \bar{B}) = \bar{B}$; we define also $p: X_2' \to X_2$ by $p(B, \bar{B}) = B$. We have a cartesian diagram
\begin{equation}\label{eq:1.6.2}
\begin{array}{ccc}
X_2' & \xrightarrow{\delta'} & X(sw) \\
\downarrow{p} & & \downarrow{F} \\
X_2 & \xrightarrow{\delta} & X(F(sw))
\end{array}
\end{equation}
For any $\bar{B} \in X(sw)$ we denote by $s\bar{B}$ the unique Borel subgroup such that $(\bar{B}, s\bar{B}) \in O(s)$, $(s\bar{B}, F\bar{B}) \in O(w)$. It is easy to see that $\delta'^{-1}(\bar{B})$ can be identified with the set of Borel subgroups $B$ such that $(B, \bar{B}) \in O(s)$, $B \neq s\bar{B}$; it follows that $\delta'^{-1}(\bar{B})$ is an affine line with a point removed. Thus $X_2'$ is (via $\delta'$) a line bundle over $X(sw)$ with the zero-section removed. Since the vertical arrows in \eqref{eq:1.6.2} induce equivalences of \'etale sites, we have a canonical exact sequence (see \cite{SGA4.5}):
\begin{equation}\label{eq:1.6.3}
\cdots \longrightarrow H_c^{i-2}(X(sw)) \xrightarrow{\alpha} H_c^i(X_2) \longrightarrow H_c^{i-2}(X(sw))(-1) \longrightarrow H_c^{i-1}(X(sw)) \xrightarrow{\alpha} \cdots
\end{equation}
It can be proved that the maps $\alpha$ in \eqref{eq:1.6.3} are zero; this fact will not be used here.

Note that \eqref{eq:1.6.1} and \eqref{eq:1.6.3} are $G^F$-equivariant and that $G^F$ acts trivially on $\overline{\mathbb{Q}}_\ell(-1)$. It follows that for any $g \in G^F$:
\[
\tr(g^*, H_c^*(X_2)) = 0, \qquad \tr(g^*, H_c^*(X_1)) = \tr(g^*, H_c^*(X(w)).
\]
If one uses the exact sequence
\[
\cdots \longrightarrow H_c^{i-1}(X_1) \longrightarrow H_c^i(X_2) \longrightarrow H_c^i(X(w')) \longrightarrow H_c^i(X_1) \longrightarrow \cdots
\]
it follows that
\[
\tr(g^*, H_c^*(X(w'))) = \tr(g^*, H_c^*(X_1)) + \tr(g^*, H_c^*(X_2)) = \tr(g^*, H_c^*(X(w)))
\]
and 1.6 is proved in this case.

\textit{The general case.} It suffices to treat the case where $w' = swF(s)$ for a fundamental reflection $s$. By permuting, if necessary, $w$ and $w'$ ($w = sw'F(s)$), we may even assume that $l(w') \ge l(w)$. If $l(w') > l(w)$ we are in case 2. If $l(w') = l(w)$, the following lemma shows that either we are in case 1 (with $w$ and $w'$ possibly interchanged) or that $w = w'$.

\begin{lemma}[1.6.4]
Let $s$, $t$ be two fundamental reflections in $W$ and let $w \in W$ be such that $l(w) = l(swt)$. Then either $w = swt$, or
\[
l(sw) = l(w) - 1, \qquad \text{or } l(wt) = l(w) - 1.
\]
\end{lemma}

Let $w = s_1 s_2 \cdots s_k$ be a reduced expression for $w$ (0.4). Assume that $l(wt) = l(w) + 1$; then $wt = s_1 s_2 \cdots s_k t$ is also a reduced expression. We have $l(swt) = l(wt) - 1$. It follows (cf.\ Bourbaki, \cite[Ch.\ IV, \S 1, Lemme 3]{Bou}) that either there exists $j$, $1 < j < k$ with $ss_1 \cdots s_j = s_1 \cdots s_{j-1}s_j$ or we have $ss_1 \cdots s_k = s_1 \cdots s_k t$. In the first case, we have $w = ss_1 \cdots s_{j-1}s_{j+1} \cdots s_k$ and $l(sw) = l(w) - 1$; in the second case we have $w = swt$ and the lemma is proved.

\subsection{1.7} Let us choose a maximal torus $T^*$ in $G$ and a Borel subgroup $B^* \subset G$ containing $T^*$, with unipotent radical $U^*$. The quotient $\dot{E} = G/U^*$ is a $T^*$-torsor ($=$ right principal homogeneous space of $T^*$) over $X = G/B^*$. For $x \in X$, the fibre $\dot{E}(x)$ of the projection $\dot{E} \to X$ is
\[
\dot{E}(x) = \{g \in G \mid ge^* = x\}/U^*
\]
where $e^*$ is the point of $X$ corresponding to $B^*$.

Let $\dot{w} \in N(T^*)$ define the element $w$ in the Weyl group $W$ via the isomorphism $\alpha(T^*, B^*): W \xrightarrow{\sim} N(T^*)/T^*$. If $x, y \in X$ are in relative position $w$, the $g$'s in $G$ such that $ge^* = x$ and $g\dot{w}e^* = y$ form a torsor $A(x, y)$ under $B^* \cap \dot{w}B^*\dot{w}^{-1} = T^* \cdot (U^* \cap \dot{w}U^*\dot{w}^{-1})$. For $g \in A(x, y)$, the class of $g\dot{w}$ in $\dot{E}(y)$ depends only on the class of $g$ in $\dot{E}(x)$. This defines a map $\dot{E}(x) \to \dot{E}(y)$, which we denote as right multiplication by $\dot{w}$. We have the formulas
\begin{align}
\label{eq:1.7.1} (ut)\dot{w} &= (u\dot{w}) \cdot \ad\dot{w}^{-1}(t), \\
\label{eq:1.7.2} u(\dot{w}\dot{v}) &= (u\dot{w})\dot{v}.
\end{align}
We will express \eqref{eq:1.7.1} by saying that $\dot{w}$ is a $w$-map of $T^*$-torsors. It is induced by a $w$-map of $T^*$-torsors over $O(w)$
\[
\dot{w}: \pr_2^* \dot{E} \longrightarrow \pr_1^* \dot{E}.
\]
Assume that $w = w_1 w_2$, $\dot{w} = \dot{w}_1\dot{w}_2$ and that $l(w) = l(w_1) + l(w_2)$; then for $x, y, z \in X$ with $(x, y) \in O(w_1)$ and $(y, z) \in O(w_2)$, we have $(x, z) \in O(w)$ and
\begin{equation}\label{eq:1.7.3}
u \dot{w} = (\nu \dot{w}_1)\dot{w}_2.
\end{equation}

\subsection{1.8} We now assume that $T^*$ and $B^*$ are $F$-stable. The identification $\alpha(T^*, B^*)$ of $T^*$ and $N(T^*)/T^*$ with the torus $T$ and the Weyl group $W$ is then compatible with $F$. For $w$ in the Weyl group, we denote by $T(w)$ the torus $T^*$, provided with the rational structure for which the Frobenius is $\ad(\dot{w}) \circ F$. We have
\[
T(w)^F = \{t \in T^* \mid \ad(\dot{w})F(t) = t\}.
\]
For $x \in X$, the Frobenius map induces a map $F: \dot{E}(x) \to \dot{E}(F(x))$, with $F(ut) = F(u)F(t)$. For $x \in X(w)$, we put
\[
\dot{E}(x, \dot{w}) = \{u \in \dot{E}(x) \mid F(u) = u\dot{w}\}.
\]
This is a $T(w)^F$-torsor. The $\dot{E}(x, \dot{w})$ are the fibres of a map
\[
\dot{w}: \dot{X}(\dot{w}) \longrightarrow X(w),
\]
with $\dot{X}(\dot{w}) \subset \dot{E}$ a $T(w)^F$-torsor over $X(w)$. The action of $G$ on $\dot{E}$ restricts to an action of $G^F$ on $\dot{X}(\dot{w})$.

Up to isomorphism, the $G^F$-equivariant $T(w)^F$-torsor $\dot{X}(\dot{w})$ over $X(w)$ is independent of the lifting $\dot{w}$ of $w$ in $N(T^*)$: for $\dot{w}' = \dot{w}t$, there exists $t_1$ with $\ad w^{-1}(t_1)F(t_1)^{-1} = t$ and the map $u \mapsto ut_1$ induces an isomorphism $\dot{X}(\dot{w}) \xrightarrow{\sim} \dot{X}(\dot{w}')$.

The groups $G^F$ and $T(w)^F$ act on $H_c^*(\dot{X}(\dot{w}), \overline{\mathbb{Q}}_\ell)$ by transport of structure. For any $\theta \in \Hom(T(w)^F, \overline{\mathbb{Q}}_\ell^*)$ we denote by $H_c^*(\dot{X}(\dot{w}), \overline{\mathbb{Q}}_\ell)_\theta$ the subspace of $H_c^*(\dot{X}(\dot{w}), \overline{\mathbb{Q}}_\ell)$ on which $T(w)^F$ acts by $\theta$.

\begin{definition}[1.9]
$R^\theta(w)$ is the virtual representation
\[
\sum_i (-1)^i H_c^i(\dot{X}(\dot{w}), \overline{\mathbb{Q}}_\ell)_\theta
\]
of $G^F$ (an element of the Grothendieck group of representations of $G^F$ over $\overline{\mathbb{Q}}_\ell$).
\end{definition}

The character $\theta$ can be used to transform the $T(w)^F$-torsor $\dot{X}(\dot{w})$ into a local system of $\overline{\mathbb{Q}}_\ell$-vector spaces of rank one $\mathcal{F}_\theta$ over $X(w)$, provided with $\theta: \dot{X}(\dot{w}) \to \mathcal{F}_\theta^\times$, $(x, t) = \theta(x)\theta(t)$.

The morphism $\dot{w}: \dot{X}(\dot{w}) \to X(w)$ is finite and
\[
\dot{w}_* \overline{\mathbb{Q}}_\ell = \bigoplus_\theta \mathcal{F}_\theta.
\]
The sheaf $\mathcal{F}_\theta$ is the subsheaf of $\dot{w}_* \overline{\mathbb{Q}}_\ell$ on which $T^F$ acts by $\theta$, hence
\[
H_c^*(\dot{X}(\dot{w}), \overline{\mathbb{Q}}_\ell)_\theta = H_c^*(X(w), \mathcal{F}_\theta).
\]
In particular, for $\theta = 1$,
\[
R^1(w) = \sum_i (-1)^i H_c^i(X(w), \overline{\mathbb{Q}}_\ell)
\]
so that Definition 1.9 is compatible with (1.5).

\begin{example}[1.10]
For $w = \dot{w} = e$, $\dot{w}: \dot{X}(\dot{w}) \to X(w)$ becomes the projection $\dot{w}: G^F/U^{*F} \to G^F/B^{*F}$, and $R^\theta(w)$ is the representation of $G^F$ on the space of functions on $G^F$ satisfying
\[
f(gtu) = \theta(t)^{-1}f(g)
\]
$G^F$ acting by $(g \cdot f)(x) = f(g^{-1}x)$ (induced representation).
\end{example}

\subsection{1.11} The Borel subgroup $\ad g B^*$ is in $X(w)$ if and only if $\ad g B^*$ and $\ad Fg B^*$ are in relative position $w$, i.e., if and only if $g^{-1}Fg \in B^*\dot{w}B^*$ (where $\dot{w} \in N(T^*)$ represents $w$):
\begin{equation}\label{eq:1.11.1}
X(w) = \{g \in G \mid g^{-1}Fg \in B^*\dot{w}B^*\}/B^*.
\end{equation}
If a Borel subgroup $B$ is in $X(w)$, one can find $g \in G$ such that $\ad g B^*$ is $B$ and $\ad g \ad\dot{w} B^* = FB$:
\begin{equation}\label{eq:1.11.2}
X(w) = \{g \in G \mid g^{-1}Fg \in \dot{w}B^*\}/B^* \cap \ad\dot{w}B^*
\end{equation}
where $B^* \cap \ad\dot{w}B^* = T^* \cdot (U^* \cap \ad\dot{w}U^*)$. Changing $g$ to $gt$, we can normalize $g$ so that we have also $g^{-1}Fg \in \dot{w}U^*$:
\begin{equation}\label{eq:1.11.3}
X(w) = \{g \in G \mid g^{-1}Fg \in \dot{w}U^*\}/T(w)^F \cdot (U^* \cap \ad\dot{w}U^*).
\end{equation}
A point in $\dot{X}(\dot{w})$ is defined by a Borel subgroup $B$, plus $g \in G$ such that $\ad g B^* = B$, $\ad g \ad\dot{w} B^* = FB$ and $g\dot{w} = Fg \mod U^*$:
\begin{equation}\label{eq:1.11.4}
\dot{X}(\dot{w}) = \{g \in G \mid g^{-1}Fg \in \dot{w}U^*\}/U^* \cap \ad\dot{w}U^*.
\end{equation}

\begin{corollary}[1.12]
The following assertions are equivalent:
\begin{enumerate}
\item[(i)] $X(w)$ is affine;
\item[(ii)] $\dot{X}(\dot{w})$ is affine;
\item[(iii)] Let $\rho$ be the action of $U^* \cap \ad\dot{w}U^*$ on $U^*$ defined by
\[
\rho(u)v = \ad\dot{w}^{-1}(u) v F(u^{-1});
\]
then $U^*/\rho(U^* \cap \ad\dot{w}U^*)$ is affine.
\end{enumerate}
\end{corollary}

Put $S = \{g \in G \mid g^{-1}Fg \in \dot{w}U^*\} \cdot \dot{w}$. The map $f: S \to U^*: g \mapsto \dot{w}^{-1}g^{-1}Fg$ induces an isomorphism $G^F\backslash S \xrightarrow{\sim} U^*$ and is such that for $u \in U^* \cap \ad\dot{w}U^*$,
$f(gu) = \rho(u)^{-1}f(g)$. Hence,
\[
G^F\backslash \dot{X}(\dot{w}) \xrightarrow{\sim} U^*/\rho(U^* \cap \ad\dot{w}U^*).
\]
As $\dot{X}(\dot{w})/T(w)^F = X(w)$, it only remains to use the fact that a space and a quotient of it by a finite group are simultaneously affine or not.

We will have to use another description of $\dot{w}: \dot{X}(\dot{w}) \to X(w)$. First, an easy lemma:

\begin{lemma}[1.13]
Let $\mathcal{J}$ be the set of pairs $(T, B)$, $T$ an $F$-stable maximal torus and $B$ a Borel subgroup containing $T$. The map $h$ which to $(T, B)$ associates the relative position of $B$ and $FB$ induces a bijection
\[
G^F\backslash \mathcal{J} \xrightarrow{\sim} W.
\]
\end{lemma}

The proof will be given in (1.15).

If we use $(T, B) \in \mathcal{J}$ to identify $W$ with $N(T)/T$, we have
\begin{equation}\label{eq:1.13.1}
h(T, \ad w_1^{-1} B) = w_1^{-1} h(T, B) F(w_1),
\end{equation}
where $w_1 \in N(T)$ represents $w$, hence

\begin{corollary}[1.14]
The map $h$ induces a bijection
\[
\{G^F\text{-conjugacy classes of }F\text{-stable maximal tori}\} \xrightarrow{\sim} W^F.
\]
\end{corollary}

Here is another description of $h$: for $(T, B) \in \mathcal{J}$, if $\alpha: T \to T$ and $\sigma: W \to W$ ($W = N(T)/T$) are defined by $(T, B)$, then
\[
\alpha F \alpha^{-1} = \ad h(T, B) \circ F: T \longrightarrow T.
\]
To give $h(T, B)$ is the same as to give $\ad h(T, B) \circ F \in W \circ F \subset \End(T)$, and to give the $F$-conjugacy class of $h(T, B)$ is the same as to give $\ad h(T, B) \circ F$ up to $W$-conjugacy.

\subsection{1.15} The space of maximal tori of $G$ can be identified with the homogeneous space $G/N(T^*)$, and the space of maximal tori marked by a containing Borel subgroup can be identified with $G/T^*$. The group $T^*$ being the connected component of $N(T^*)$, (1.13) is a special case of the general result described below.

Let $G_0$ be a connected algebraic group over $\mathbb{F}_q$ and let $\pi: E_0 \to E_0$ be a morphism of $G_0$-homogeneous spaces. We denote by $G$, $E$, $E$ the corresponding objects over $k$, and we assume that the stabilizer $S(e)$ of $e \in E$ is the connected component of the stabilizer $S(\pi(e))$ of $\pi(e) \in E$. Since any $G_0$-homogeneous space has a rational point, the existence of $E_0$ imposes no condition on $E_0$.

The groups $S(e)/S(e)^\circ$ form a local system on $E$, which becomes constant on $E$; we denote by $W$ its constant value on $E$. For $e \in E$, we have an isomorphism $\alpha_\pi(e): W \xrightarrow{\sim} S(\pi(e))/S(e)$ and, for $f = ge$, we have $\alpha_\pi(f) = \ad g \circ \alpha_\pi(e)$. We let the group $W$ act on $E$ on the right, by $e \cdot w = \alpha_\pi(e)(w) \cdot e$; in this way, $E$ becomes a $W$-torsor ($=$ principal homogeneous space) over $E$.

The group $W$ is acted on by $F$, with $F(e \cdot w) = F(e) \cdot F(w)$. The set $W^F$ of $F$-conjugacy classes in $W$ is the set of orbits of the action of $W$ on itself by $w \mapsto w^{-1}F(w)$.

\begin{proposition}[1.16]
For $e \in E^F$ and $\tilde{e} \in E$ above it, define $h(e, \tilde{e}) \in W$ by
\[
F(\tilde{e}) = \tilde{e} \cdot h(e, \tilde{e}).
\]
\begin{enumerate}
\item[(i)] The map $h$ induces a bijection from the set of $G^F$-orbits in $\{(e, \tilde{e}) \mid e \in E^F, \pi(\tilde{e}) = e\}$ to $W$.
\item[(ii)] We have $h(e, \tilde{e}w) = w^{-1}h(e, \tilde{e})F(w)$. Hence the map $h$ induces a bijection from $G^F\backslash E^F$ to the set of $F$-conjugacy classes in $W$.
\end{enumerate}
\end{proposition}

We will only prove (i). Let $Y$ be the set of $\tilde{e} \in E$ such that $\pi(\tilde{e}) \in E^F$. If $e_0 \in E^F$, the map $g \mapsto g \cdot e_0$ identifies $X$ with $\{g \in G \mid g^{-1}Fg \in S(\pi(e_0))\}/S(e_0)$. The Lang isogeny $g \mapsto g^{-1}Fg$ is an isomorphism $G^F\backslash G \to G$, hence the map $g \cdot e_0 \mapsto g^{-1}Fg$ induces a bijection
\[
G^F\backslash X \xrightarrow{\sim} S(\pi(e_0))/(S(e_0) \text{ acting by } s \mapsto s^{-1}xF(s))
\]
The orbits of this action of $S(e_0)$ are just the usual cosets, and the resulting bijection $G^F\backslash X \xrightarrow{\sim} W$ is the $h$ above.

\begin{definition}[1.17]
Let $T$ be an $F$-stable maximal torus and let $B$ be a Borel subgroup containing $T$, with unipotent radical $U$; let $w$ be the relative position of $B$ and $FB$.
\begin{enumerate}
\item[(i)] $X_T \subset B$ is $X(w)$. The map $g \mapsto \ad g B$ induces isomorphisms
\begin{align*}
X_T \subset B &\cong \{g \in G \mid g^{-1}Fg \in BF(B)\}/B \\
&\cong \{g \in G \mid g^{-1}Fg \in FU\}/T^F \cdot (U \cap FU).
\end{align*}
\item[(ii)] $\dot{X}_T \subset B$ is $\{g \in G \mid g^{-1}Fg \in FU\}/U \cap FU$.
\end{enumerate}
We have a projection map $\dot{w}: \dot{X}_T \subset B \to X_T \subset B$, for which $\dot{X}_T \subset B$ is a $G^F$-equivariant $T^F$-torsor over $X_T \subset B$; $G^F$ acts by left multiplication and $T^F$ by right multiplication (it normalizes $U \cap FU$).
\end{definition}

\subsection{1.18} Let $\dot{w} \in N(T^*)$ be a representative of $w$. If $x' \in G$ is such that $\ad x'(T^*, B^*) = (T, B)$ then $B$ and $FB = \ad Fx' B^*$ are in relative position $w$, and $FB$ contains $T$, hence $FB = \ad x' \ad\dot{w} B^*$ and $x'^{-1}F(x') \in \dot{w}B^* \cap N(T^*) = \dot{w}T^*$. Replacing $x'$ by $x = x't$ ($t \in T^*$) one can achieve $x^{-1}F(x) = \dot{w}$. The $x$ such that $\ad x(T^*, B^*) = (T, B)$ and $x^{-1}F(x) = \dot{w}$ form a $T^{*F}$-torsor. For such an $x$, $\ad x$ induces an isomorphism $T(w) \xrightarrow{\sim} T$ (hence $T(w)^F \xrightarrow{\sim} T^F$); this isomorphism is independent of $x$.

\begin{proposition}[1.19]
Let $T$, $B$, $U$, $w$ be as in (1.17) and $x$, $\dot{w}$ as in (1.18). The map $g \mapsto gx^{-1}$ induces an isomorphism from the $G^F$-equivariant $T(w)^F$-torsor $\dot{X}(\dot{w})$ over $X(w)$ (or rather its model \eqref{eq:1.11.4}) to the $G^F$-equivariant $T^F$ torsor $\dot{X}_T \subset B$ over $X_T \subset B$ (or rather its model (1.17)).
\end{proposition}

This is a straightforward computation.

\subsection{1.20} The cohomology of $\dot{X}_T \subset B$ is acted on by $G^F$ and $T^F$. For any character $\theta: T^F \to \overline{\mathbb{Q}}_\ell^*$ we put as in (1.8):
\[
R_T^\theta \subset B = \sum_i (-1)^i H_c^i(\dot{X}_T \subset B, \overline{\mathbb{Q}}_\ell)_\theta
\]
(an element in the Grothendieck group of representations of $G^F$ over $\overline{\mathbb{Q}}_\ell$).

By (1.19), for $x$ as in (1.18), we have
\[
R_T^\theta \subset B = R^{\theta \circ \ad x}(w).
\]
We shall see in Chapter 4 that $R_T^\theta \subset B$ is independent of $B$. For $g \in G^F$ such that $\ad g$ carries $T$, $B$, and $\theta$ to $T'$, $B'$, and $\theta'$, we have clearly $R_T^\theta \subset B = R_{T'}^{\theta'} \subset B'$; hence $R_T^\theta \subset B$ will eventually depend only on the $G^F$-conjugacy class of $T$ and on the orbit of $\theta$ under $(N(T)/T)^F$.

The end of this chapter will be used in the proof of 7.10 only.

\subsection{1.21} \textit{Isogenies.} Let $T$ be an $F$-stable maximal torus of $G$, and let $Z$ be the centre of $G$. We denote $G \times_Z T$ the quotient of $G \times T$ by the subgroup $\{(z, z^{-1}) \mid z \in Z\}$.

Let $B$ be a Borel subgroup containing $T$. The action $x \mapsto gxt$ of $G^F \times T^F$ on $\dot{X}_T \subset B$ is induced by an action of $(G \times_Z T)^F$, given by the same formula.

\begin{corollary}[1.22]
On $H_c^*(\dot{X}_T \subset B, \overline{\mathbb{Q}}_\ell)_\theta$, $Z^F$ acts by the character $\theta \mid Z^F$. In particular, on any irreducible representation occurring in $R_T^\theta \subset B$, $Z^F$ acts by $\theta \mid Z^F$.
\end{corollary}

Let $\pi: \tilde{G} \to G$ be the simply connected covering of the derived group of $G$. $\tilde{T} = \pi^{-1}(T)$, $\tilde{B} = \pi^{-1}(B)$ and let $\tilde{Z}$ be the centre of $\tilde{G}$.

\begin{proposition}[1.23]
One has $T^F/\pi(\tilde{T}^F) \xrightarrow{\sim} G^F/\pi(\tilde{G}^F)$.
\end{proposition}

Injectivity is clear; we have to check the surjectivity of the map $(\tilde{T} \times \tilde{G})^F \to G^F$ induced by $\varphi': \tilde{T} \times \tilde{G} \to G: (t, g) \mapsto t\pi(g)$. Via $\varphi'$, $\tilde{T} \times \tilde{G}$ is a $\tilde{T}$-torsor over $G$ (with $(t, g) * \tilde{t} = (t\tilde{t}, \tilde{t}^{-1}g)$), and one applies Lang's theorem to the connected group $\tilde{T}$.

\subsection{1.24} Let $h: A \to B$ be a homomorphism of finite groups and let $X$ be a space on which $A$ acts. The induced space $\Ind_A^B(X)$ (unique up to unique isomorphism) is any $B$-space $I$, provided with an $A$-equivariant map $\lambda: X \to I$, such that for any $B$-space $Y$, $\Hom_B(I, Y) \xrightarrow{\sim} \Hom_A(X, Y)$. One has $\Ind_A^B(X) = \coprod_{b \in B/hA} b \cdot \lambda(X)$, and $\lambda(X) \xrightarrow{\sim} \Ker(h) \backslash X$.

\begin{proposition}[1.25]
The $(G \times_Z T)^F$-space $\dot{X}_T \subset B$ is induced by the $(\tilde{G} \times_{\tilde{Z}} \tilde{T})^F$ space $\dot{X}_{\tilde{T}} \subset \tilde{B}$.
\end{proposition}

\begin{lemma}[1.26]
In the diagram
\[
\begin{array}{ccccccccccc}
0 & \longrightarrow & \Ker & \longrightarrow & \tilde{T}^F & \longrightarrow & T^F & \longrightarrow & \Coker & \longrightarrow & 0 \\
& & \downarrow{(1)} & & \downarrow{(t,1)} & & \downarrow{(t,1)} & & \downarrow{(2)} & & \\
0 & \longrightarrow & \Ker & \longrightarrow & (\tilde{T} \times_{\tilde{Z}} \tilde{G})^F & \longrightarrow & (T \times_Z G)^F & \longrightarrow & \Coker & \longrightarrow & 0 \\
& & & & & & \downarrow & & & \\
& & & & & & (G/Z)^F & = & (G/Z)^F &
\end{array}
\]
the maps (1) and (2) are isomorphisms.
\end{lemma}

Indeed, $\tilde{T} \times_{\tilde{Z}} \tilde{G}$ (resp.\ $T \times_Z G$) is a $\tilde{T}$ (resp.\ $T$)-torsor over $G/Z = G/Z$, hence, since $\tilde{T}$ and $T$ are connected, $(\tilde{T} \times_{\tilde{Z}} \tilde{G})^F$ is a $\tilde{T}^F$-torsor over $(G/Z)^F$, and $(T \times_Z G)^F$ is the induced $T^F$-torsor.

\textit{Proof of 1.25.} By 1.26, we are reduced to prove that $\dot{X}_T \subset B$, as a $T^F$-space, is induced by the $\tilde{T}^F$-space $\dot{X}_{\tilde{T}} \subset \tilde{B}$. The spaces $\dot{X}_{\tilde{T}} \subset \tilde{B}$ and $\dot{X}_T \subset B$ are indeed respectively $\tilde{T}^F$ and $T^F$-torsor over $X_T \subset B = X_{\tilde{T}} \subset \tilde{B}$ and $X_T \subset B$ is hence the $T^F$-torsor induced by the $\tilde{T}^F$-torsor $\dot{X}_{\tilde{T}} \subset \tilde{B}$.

\begin{corollary}[1.27]
Let $\theta$ be a character of $G^F/\pi(\tilde{G}^F)$. We denote again by $\theta$ its restriction to $T^F$. One has $R_T^\theta \subset B = \theta \cdot R_T^1 \subset B$.
\end{corollary}

It follows from 1.25 that
\[
H_c^*(\dot{X}_T \subset B, \overline{\mathbb{Q}}_\ell) = \Ind_{(\tilde{G} \times_{\tilde{Z}} \tilde{T})^F}^{(G \times_Z T)^F} H_c^*(\dot{X}_{\tilde{T}} \subset \tilde{B}, \overline{\mathbb{Q}}_\ell).
\]
The character $\theta(gt)$ of $(G \times_Z T)^F$ is trivial on the image of $\tilde{G} \times_{\tilde{Z}} \tilde{T}$. The induced representation we consider is hence isomorphic to its tensor product with $\theta(gt)$, and 1.27 is a formal consequence of this.

\section{Examples in the classical groups}

\subsection{2.1} Let $V$ be an $n$-dimensional vector space over $k$ and put $G = GL(V)$. If $b = (b_1, \dots, b_n)$ is a basis of $V$, we may take for $T^*$ the group of diagonal matrices and for $B^*$ the group of upper triangular matrices. The Weyl group lifts into the subgroup of $N(T^*)$ consisting of the $\dot{w}$'s inducing a permutation of basis vectors. In this case, $T$, $W$ (1.1), $X$ (1.2), $\dot{E}$, $\dot{w}$ (1.7), have the following alternative description.

(a) $T = \mathbb{G}_m^n$, $W = \mathfrak{S}_n$, the fundamental reflections are the transpositions $(i, i+1)$ and the action of $W$ on $T$ is by permutation.

(b) $X$ is the space of complete flags $D_1 \subset \cdots \subset D_{n-1}$ in $V$: a flag $D$ is an increasing filtration of $V$ with $\dim D_i = i$ for $1 \le i \le n-1$.

(c) $\dot{E}$ is the space of complete flags marked by non-zero vectors $e_i \in D_i/D_{i-1} = \Gr_i^D(V)$ ($1 \le i \le n$), where we use the convention $D_0 = 0$, $D_n = V$; $T$ acts on $\dot{E}$ by $(D, (e_i))(\lambda_i) = (D, (\lambda_i e_i))$. This is a $G$-equivariant $T$-torsor over $X$.

(d) If $D'$ and $D''$ are two flags, their relative position is labelled by the permutation $w$ such that $\Gr_{w(i)}^{D'} \cap \Gr_i^{D''}(V) \neq 0$. The isomorphisms
\[
\Gr_{w(i)}^{D'}(V) \xleftarrow{\sim} \Gr_{w(i)}^{D'} \cap \Gr_i^{D''}(V) \xrightarrow{\sim} \Gr_i^{D''}(V)
\]
induce a $w$-isomorphism between the $T$-torsor $\dot{E}(D')$ of markings of $D'$ and $\dot{E}(D'')$: $e \mapsto e \circ w$. When $w$ is the $n$-cycle $(1, \dots, n)$, $D'$ and $D''$ are in relative position $w$ if and only if
\[
D_i + D_i' = D_{i+1}' \ (1 \le i \le n-1) \quad \text{and} \quad D_1 + D_1' = V.
\]

\subsection{2.2} We now take $k$ and $\mathbb{F}_q$ as in (0.1.1) and assume that $V$ is provided with an $\mathbb{F}_q$-structure. Frobenius maps are then defined. For $w = (1, \dots, n)$, the condition for a flag $D$ to be in relative position $w$ with its image $FD$ by Frobenius is that $D$ be the flag
\[
D_1 \subset D_1 + FD_1 \subset D_1 + FD_1 + F^2D_1 \subset \cdots \quad \text{and that } V = \bigoplus_0^{n-1} F^i D_1.
\]
If we denote by $P(V)$ the set of homogeneous lines in $V$, the map $D \mapsto D_1$ is an isomorphism from $X(w)$ to the set of all $x \in P(V)$ which do not lie on any $\mathbb{F}_q$-rational hyperplane. A marking $e$ of $F$ is such that $F(e) = e \circ w$ if and only if
\[
e_2 \equiv F(e_1) \pmod{e_1}, \ e_3 \equiv F^2(e_1) \pmod{e_1, F(e_1)}, \dots, e_n \equiv F^{n-1}(e_1) \pmod{e_1, F(e_1), \dots, F^{n-2}(e_1)}
\]
and
\[
e_1 \equiv F^n(e_1) \pmod{F(e_1), \dots, F^{n-1}(e_1)}.
\]
$e$ is defined by $e_1 \in D_1$ subject to the condition that
\[
e_1 \wedge F(e_1) \wedge \cdots \wedge F^{n-1}(e_1) = F^n(e_1) \wedge F(e_1) \wedge \cdots \wedge F^{n-1}(e_1),
\]
i.e.:
\begin{equation}\label{eq:2.2.1}
F(e_1 \wedge \cdots \wedge F^{n-1}(e_1)) = (-1)^{n-1}(e_1 \wedge \cdots \wedge F^{n-1}(e_1)).
\end{equation}
If $(x_i)$ are the coordinates of $e_1$ with respect to some rational basis, the condition \eqref{eq:2.2.1} can be rewritten
\begin{equation}\label{eq:2.2.2}
(-1)^{n-1}(\det(x_i^{q^j})_{1 \le i,j \le n})^{q-1} = 1.
\end{equation}
The form on the left is invariant under $GL(n, \mathbb{F}_q)$. Up to a scalar factor, it is the product of all non-zero $\mathbb{F}_q$-rational linear forms. The map $(D, e) \mapsto e_1$ induces an isomorphism of $\dot{X}(\dot{w})$ with the affine hypersurface \eqref{eq:2.2.2}. This hypersurface is stable under $x \mapsto \lambda x$ for $\lambda \in \mathbb{F}_q^*$, and this is the action of $T(w)^F$.

Our work has been inspired by results of Drinfeld, who proved that the discrete series representations of $SL(2, \mathbb{F}_q)$ occur in the cohomology of the affine curve $xy^q - x^qy = 1$ (the form $xy^q - x^qy$ is $SL_2(\mathbb{F}_q)$-invariant).

Let $\varphi(q, r, n)$ be the number of $\mathbb{F}_{q^r}$-rational points of $X(w)$, $r \ge 1$. We have the following

\begin{proposition}[2.3]
\[
\varphi(q, r, n) = \prod_{i=1}^{n-1}(q^r - q^i).
\]
\end{proposition}

\textit{Proof.} We define a partition $P(V) = X_0 \cup X_1 \cup \cdots \cup X_n$, as follows: $X_i$ is the set of points $x \in P(V)$ such that $x, F(x), F^2(x), \dots$ span a linear subspace of dimension $i$ of $P(V)$. This partition is invariant under Frobenius, and clearly $X_i$ has precisely
\[
q^{\frac{i(i-1)}{2}}\binom{n}{i}_{q^{-1}}(1-q^{n-i})\cdots(1-q^{n-1})
\]
$\mathbb{F}_q$-rational points. It follows that
\[
\varphi(q, r, n) = -\sum_{i=0}^{n-1} \varphi(q, r, i+1)(1-q^{n-i})\cdots(1-q^{n-1}).
\]
This shows that, for fixed $n$, $\varphi(q, r, n)$ is a polynomial of degree $(n-1)$ in $Q = q^r$, with coefficient polynomials in $q$, and leading term $Q^{n-1}$. Since $\varphi(q, r, n) = 0$ for $1 \le r < n-1$, this polynomial must be divisible by $(Q-q)$, $(Q-q^2), \dots, (Q-q^{n-1})$. It follows that
\[
\varphi(q, r, n) = (Q-q)(Q-q^2)\cdots(Q-q^{n-1})
\]
and the proposition is proved.

\subsection{2.4} Let $V$ be a vector space as in (2.2), with a fixed non-degenerate symplectic form $\langle, \rangle$ defined over $\mathbb{F}_q$. We must have $n = 2m$. The symplectic group $Sp(V)$ is then a group as in (0.1.1). Let $Y$ be the set of all complete isotropic flags $D_1 \subset D_2 \subset \cdots \subset D_m$ in $V$ ($\dim D_i = i$) such that
\[
D_1 \neq FD_1 \subset D_2, \quad D_2 \neq FD_2 \subset D_3, \dots, D_{m-1} \neq FD_{m-1} \subset D_m, \quad D_m \neq FD_m.
\]
Then $Y$ can be identified with $X(w)$, where $w$ is a Coxeter element in the Weyl group $W$ of $Sp(V)$; if we identify $W$ with the group of all permutations $\sigma$ of $\{-m, \dots, -2, -1, 1, 2, \dots, m\}$ such that $\sigma(-i) = -\sigma(i)$ for all $i$, then $w$ is the permutation
\[
i \mapsto \begin{cases}
i+1 & (2 \le i \le m), \\
-1 \mapsto -m, & \\
i-1 & (2 \le |i| \le m), \\
1 \mapsto -m.
\end{cases}
\]
On the other hand, $Y$ (hence also $X(w)$) can be identified with the set of $x \in P(V)$ such that
\[
\langle x, F(x) \rangle = \langle x, F^2(x) \rangle = \cdots = \langle x, F^{m-1}(x) \rangle = 0, \quad \langle x, F^m(x) \rangle \neq 0.
\]
For example, if $\dim V = 4$, this is just the set of all $x \in P(V)$ such that $\langle x, F(x) \rangle = 0$, $\langle x, F^2(x) \rangle \neq 0$. Note that the equation $\langle x, F(x) \rangle = 0$ defines a non-singular surface $S$ in $P(V)$ and that $\langle x, F^2(x) \rangle \neq 0$ means that we remove from $S$ a union of rational curves, one for each isotropic plane in $V$, defined over $\mathbb{F}_q$. One can prove that in this case ($n = 4$), the number of $\mathbb{F}_{q^r}$-rational points of $X(w)$ is given by:
\[
q^{2r} - q \cdot (1+q)q^{2r} + q(1-q)^2(-q)^r + q^3 \cdot (q^r - q)(q^r - 3) \quad \text{($r$ odd)}.
\]

\section{A fixed point formula}

\subsection{3.1} Let $X$ be a scheme, separated and of finite type over $k$ and let $\sigma: X \to X$ be an automorphism of finite order of $X$. We decompose $\sigma$ as $\sigma = s \cdot u$ where $s$ and $u$ are powers of $\sigma$ respectively of order prime to $p$ and a power of $p$. The main result of this chapter is the following:

\begin{theorem}[3.2]
With the above notations (see also 0.5),
\[
\Tr(u^*, H_c^*(X, \overline{\mathbb{Q}}_\ell)) = \Tr(u^*, H_c^*(X^s, \overline{\mathbb{Q}}_\ell)).
\]
\end{theorem}

The first step is to prove that the left hand side is an integer independent of $\ell$. One might conjecture that for any endomorphism $f$ of $X$ and each $i$, $\Tr(f^*, H_c^i(X, \overline{\mathbb{Q}}_\ell))$ is integral and independent of $\ell$, but such a more general and precise result is known only for $X$ proper and smooth (Katz and Messing, \textit{Inv.\ Math.}\ 23 (1974), 73--77).

\begin{proposition}[3.3]
With the notations of 3.1, $\Tr(\sigma^*, H_c^*(X, \overline{\mathbb{Q}}_\ell))$ is an integer independent of $\ell$ ($\ell \neq p$).
\end{proposition}

A standard specialization argument allows us to assume that $p > 0$ and that $k$ is an algebraic closure of the prime field $\mathbb{F}_p$. This is anyway the only case we will need in the rest of the paper. The scheme $X$ and $\sigma$ can then be defined over some finite extension $\mathbb{F}_q \subset k$ of $\mathbb{F}_p$; we denote by $F: X \to X$ the corresponding Frobenius endomorphism.

Let us first assume that $X$ is quasi-projective. Then, for $n \ge 1$, the composite $F \circ \sigma^n$ is the Frobenius map relative to some new way of lowering the field of definition of $X$ from $k$ to $\mathbb{F}_{q^n}$ and the Lefschetz fixed point formula for Frobenius (\cite{SGA4.5} and \cite{Gro}) shows that $\Tr((F^n\sigma)^*, H_c^*(X, \overline{\mathbb{Q}}_\ell))$ is the number of fixed points of $F^n\sigma$ ($n \ge 1$). The automorphisms $F^*$ and $\sigma^*$ of the cohomology commute; as a function of $n$, $\Tr((F^n\sigma)^*, H_c^*(X, \overline{\mathbb{Q}}_\ell))$ is hence of the form $\sum a_\lambda \lambda^n$ where $\lambda$ runs through the multiplicative group $\overline{\mathbb{Q}}_\ell^*$ of $\overline{\mathbb{Q}}_\ell$ and where $a_\lambda$ is zero for all $\lambda$ except for a finite number of them.

The functions $\mathbb{N}^+ \to \overline{\mathbb{Q}}_\ell^*: n \mapsto \lambda^n$ (for $\lambda \in \overline{\mathbb{Q}}_\ell^*$) are linearly independent. This can be checked either by using Vandermonde determinants or by appealing to Dedekind's theorem on the linear independence of characters (Bourbaki, \textit{Alg\`ebre} V, \S 7, 5). For $n \ge 1$ the numbers $\sum a_\lambda \lambda^n = |X^{F^n\sigma}|$ are rational; hence for any automorphism $\tau$ of $\overline{\mathbb{Q}}_\ell$, one has
\[
\sum_\lambda \tau(a_\lambda)\tau(\lambda)^n = \sum_\lambda \tau(a_{\tau^{-1}(\lambda)})\lambda^n = \sum_\lambda a_\lambda \lambda^n \quad (n > 1),
\]
and, by the linear independence of the functions $\lambda^n$,
\[
\tau(a_\lambda) = a_{\tau(\lambda)}.
\]
In particular, $\Tr(\sigma^*, H_c^*(X, \overline{\mathbb{Q}}_\ell)) = \sum a_\lambda$ is invariant by any automorphism of $\overline{\mathbb{Q}}_\ell$, hence rational. Similarly, as $\Tr((F^n\sigma)^*, H_c^*(X, \overline{\mathbb{Q}}_\ell)) = |X^{F^n\sigma}|$ is independent of $\ell$, for any isomorphism $\tau: \overline{\mathbb{Q}}_\ell \xrightarrow{\sim} \overline{\mathbb{Q}}_{\ell'}$ one has
\[
\tau \Tr(\sigma^*, H_c^*(X, \overline{\mathbb{Q}}_\ell)) = \Tr(\sigma^*, H_c^*(X, \overline{\mathbb{Q}}_{\ell'})),
\]
hence the rational number $\Tr(\sigma^*, H_c^*(X, \overline{\mathbb{Q}}_\ell))$ is independent of $\ell$.

Since the automorphism $\sigma$ has finite order, $\Tr(\sigma^*, H_c^*(X, \overline{\mathbb{Q}}_\ell))$ is a sum of roots of unity, hence an algebraic integer. Being rational it is an ordinary integer.

If we do not assume $X$ to be quasi-projective, we can either repeat the previous argument by working with algebraic spaces or reduce to the quasi-projective case: if $(X_i)$ is a finite partition of $X$ into locally closed quasi-projective subschemes stable under $\sigma$, then (cf.\ \cite{SGA4.5})
\begin{equation}\label{eq:3.3.1}
\Tr(\sigma^*, H_c^*(X, \overline{\mathbb{Q}}_\ell)) = \sum_i \Tr(\sigma^*, H_c^*(X_i, \overline{\mathbb{Q}}_\ell)).
\end{equation}

\subsection{3.4} \textit{Proof of 3.2.} Let $(X_i)$ be a partition of $X$ into locally closed subschemes stable under $\sigma$ and such that, on each $X_i$, $\sigma$ defines a free action of a cyclic group. By applying \eqref{eq:3.3.1} to the decompositions $X = \bigcup X_i$ and $X^s = \bigcup X_i^s$, one reduces to the case where $\sigma$ generates a free action of a cyclic group $H$. In this case, either

(a) $\sigma$ is of order a power of $p$, hence $s = \id$, $u = \sigma$ and (3.2) is trivial, or

(b) the order of $\sigma$ is divisible by a prime number $\ell' \neq p$ and we must prove that $\Tr(\sigma^*, H_c^*(X, \overline{\mathbb{Q}}_\ell)) = 0$. As
\begin{equation}\label{eq:3.3}
\Tr(\sigma^*, H_c^*(X, \overline{\mathbb{Q}}_\ell)) = \Tr(\sigma^*, H_c^*(X, \overline{\mathbb{Q}}_{\ell'})),
\end{equation}
we may as well assume that $\ell = \ell'$. Let us now grant the

\begin{proposition}[3.5]
Let $H$ be a finite group acting freely on $X$. Then $\chi(h) = \Tr(h^*, H_c^*(X, \overline{\mathbb{Q}}_\ell))$ is the character of a virtual projective $\mathbb{Z}_\ell[H]$-module.
\end{proposition}

If $H$ is abelian, and if $H'$ is its largest subgroup of order prime to $\ell$, a representation of $H$ over $\mathbb{Z}_\ell$ gives rise to a projective $\mathbb{Z}_\ell[H]$-module if and only if it is induced from a representation of $H'$. Its character vanishes then on $H - H'$ and we get the vanishing (b) by applying (3.5) to the cyclic group generated by $\sigma$.

\subsection{3.6} Let $\Lambda$ be a torsion ring with unit and let $\mathcal{F}$ be a sheaf of left $\Lambda$-modules on a scheme $Y$, with $Y$ separated and of finite type over $k$. The $\Lambda$-modules $H_c^i(Y, \mathcal{F})$ are then the cohomology modules of a finer object $R\Gamma_c(Y, \mathcal{F})$ in the derived category $D^b(\Lambda)$ of the category of $\Lambda$-modules. The following is the key to the proof of 3.5.

\begin{proposition}[3.7]
Assume $\Lambda$ to be right and left noetherian. If $\mathcal{F}$ is a constructible sheaf of projective $\Lambda$-modules, then $R\Gamma_c(Y, \mathcal{F})$ can be represented by a finite complex of projective $\Lambda$-modules of finite type.
\end{proposition}

We repeat the proof in \cite[XVII (5.2.10)]{SGA4}.

(a) The $H_c^i(Y, \mathcal{F})$ are $\Lambda$-modules of finite type, and vanish for $i > 2\dim(Y)$ (\cite[XVII (5.2.8.1) and (5.3.6)]{SGA4}).

(b) For any right $\Lambda$-module of finite type $N$, one has
\[
R\Gamma_c(Y, N \otimes_\Lambda \mathcal{F}) = N \otimes_\Lambda^L R\Gamma_c(Y, \mathcal{F})
\]
(\cite[XVII (5.2.9)]{SGA4}); the proof rests on replacing $N$ by a free resolution of $N$, to reduce to the trivial case where $N$ is free of finite type; we are allowed to use such left infinite resolutions because $R\Gamma_c$ is of finite cohomological dimension. The assumption that $N$ is of finite type is in fact unnecessary.

(c) By (a) we can represent $R\Gamma_c(Y, \mathcal{F})$ by a complex of $\Lambda$-modules $K^\bullet$, with $K^i = 0$ for $i \notin [0, 2\dim Y]$ and $K^i$ free of finite type for $i > 0$. For any right $\Lambda$-module of finite type $N$ and any $i > 0$,
\[
\Tor_i^\Lambda(N, K^0) = H^{-i}(N \otimes_\Lambda^L R\Gamma_c(Y, \mathcal{F})) \overset{(b)}{=} H^{-i} R\Gamma_c(Y, N \otimes_\Lambda \mathcal{F}) = 0.
\]
The $\Lambda$-module $K^0$ is hence flat; it is of finite type because $H^0(K^\bullet)$ is, hence it is projective.

\subsection{3.8} \textit{Proof of 3.5.} We will assume that $X$ is quasi-projective (the general case can be handled as in 3.3). Put $Y = X/H$ and denote by $\pi$ the projection $\pi: X \to Y$. The group $H$ acts on $\pi_*\mathbb{Z}/\ell^n$ and this action turns $\pi_*\mathbb{Z}/\ell^n$ into a locally constant sheaf of free $\mathbb{Z}/\ell^n[H]$-modules of rank one. By 3.7, $R\Gamma_c(Y, \pi_*\mathbb{Z}/\ell^n)$ can be represented by a complex $K_n$ of projective $\mathbb{Z}/\ell^n[H]$-modules of finite type. One has $R\Gamma_c(Y, \pi_*\mathbb{Z}/\ell^n) = R\Gamma_c(Y, \pi_*\mathbb{Z}/\ell^{n+1}) \otimes^L \mathbb{Z}/\ell^n$ (3.7 (b)), and one checks easily (\cite[XV, 3.3, Lemme]{SGA4}) that once $K_n$ is chosen, one can choose $K_{n+1}$ such that $K_n$ is the reduction mod $\ell^n$ of $K_{n+1}$. Taking a projective limit, we get a complex $K_\infty$ of projective $\mathbb{Z}_\ell[H]$-modules and an $H$-equivariant isomorphism
\[
H_c^*(X, \mathbb{Z}_\ell) = \varprojlim H_c^*(X, \mathbb{Z}/\ell^n) = \varprojlim H_c^*(Y, \pi_*\mathbb{Z}/\ell^n) = H^*(K_\infty)
\]
(the middle equality because $\pi$ is finite). We then have
\[
\chi(h) = \sum_i (-1)^i \Tr(h, K_\infty^i).
\]

\subsection{3.9} Let $X/k$ be as in (3.1), $T$ a finite abelian group of order prime to $p$, $G$ a finite group and $\rho$ an action of $T \times G$ on $X$. We assume that the action of $T$ on $X$ is free. We let $T \times G$ act on $H_c^*(X)$ by transport of structure. For any character $\theta$ of $T$ with values in $\overline{\mathbb{Q}}_\ell$, we denote by $H_{c,\theta}^*$ the subspace of $H_c^*(X)$ on which $T$ acts by $\theta$; on $H_{c,\theta}^*$, $t \in T$ acts by multiplication by $\theta(t)^{-1}$.

Let us assume for simplicity that $X$ is quasi-projective. We denote by $\pi$ the projection $\pi: X \to Y = X/T$. The group $G$ acts on $Y$. For each $g \in G$, we denote by $I(g)$ the set of connected components of the fixed point set $Y^g$ and by $Y_i^g$ the component corresponding to $i \in I(g)$.

For $y \in Y$, $\pi^{-1}(y)$ is a principal homogeneous set for the action of the abelian group $T$. If $y \in Y^g$, $g$ acts on $\pi^{-1}(y)$ and commutes with $T$, it hence acts like some $t(g, y) \in T$:
\[
gx = t(g, x) \cdot x, \quad x \in \pi^{-1}(y);
\]
$t(g, y)$ is constant for $y$ in each connected component $Y_i^g$ of $Y^g$; we denote it by $t(g, i)$.

The Theorem 3.2 will be used in the following form.

\begin{corollary}[3.10]
With the above notations, let $g = su$ be the decomposition of $g \in G$ as the product of commuting elements respectively of order prime to $p$ and a power of $p$. Then,
\[
\Tr(g^*, H_{c,\theta}^*) = \sum_{i \in I(s)} \theta(t(s, i)^{-1}) \Tr(u^*, H_c^*(Y_i^s)).
\]
\end{corollary}

For $t \in T$, the decomposition 3.1 of $\sigma = tg$ is $tg = (st) \cdot u$. We have
\[
X^{st} = \coprod_{i \in I(s)} \pi^{-1}(Y_i^s)_{t = t(s, i)}
\]
hence
\begin{align*}
\Tr(g^*, H_{c,\theta}^*) &= \frac{1}{|T|} \sum_t \theta(t)^{-1} \Tr(g^*t^*, H_c^*(X)) \\
&= \frac{1}{|T|} \sum_t \theta(t)^{-1} \sum_{i \in I(s), t = t(s, i)} \Tr(u^*, H_c^*(\pi^{-1}(Y_i^s))) \\
&= \frac{1}{|T|} \sum_{i \in I(s)} \theta(t(s, i)^{-1}) \Tr(u^*, H_c^*(\pi^{-1}(Y_i^s))).
\end{align*}
The decomposition (3.1) of the automorphism $tu$ of $\pi^{-1}(Y_i^s)$ is $t \cdot u$. For $t \neq e$, $t$ has no fixed point, hence
\[
\Tr((tu)^*, H_c^*(\pi^{-1}(Y_i^s))) = 0 \qquad (t \neq e).
\]
The endomorphism $\frac{1}{|T|}\sum_t t^*$ of $H_c^*(\pi^{-1}(Y_i^s))$ is a retraction onto the subspace $\pi^* H_c^*(Y_i^s)$, hence
\[
\frac{1}{|T|} \sum_t \Tr(u^*, H_c^*(\pi^{-1}(Y_i^s))) = \Tr(u^*, H_c^*(Y_i^s))
\]
and, substituting into the expression for $\Tr(g^*, H_{c,\theta}^*)$, we get (3.10).

\subsection{3.11} The end of this chapter will not be used in this paper. Let $G$ be a finite group acting on $X/k$ as in (3.1). For simplicity we assume $X$ to be quasi-projective. Assume further that $G$ acts freely on $X$, and put $Y = X/G$. The covering $X$ of $Y$ is said to be \textit{tame} if $Y$ can be imbedded as a Zariski open dense subset of a proper scheme $\bar{Y}$ in such a way that the $p$-Sylow-subgroups of $G$ act freely on the normalisation $\bar{X}$ of $\bar{Y}$ in $X$:
\[
\begin{array}{ccc}
\bar{X} & \longleftarrow & X \\
\downarrow & & \downarrow \\
\bar{Y} & \longleftarrow & Y.
\end{array}
\]

\begin{proposition}[3.12]
With the above notations, if $X/Y$ is tame, then the virtual representation $\sum_i (-1)^i H_c^i(X)$ of $G$ is a multiple of the regular representation.
\end{proposition}

It suffices to prove that $\Tr(g^*, H_c^*(X)) = 0$ for $g \neq e$. If $g$ is not of order a power of $p$, this follows from (3.2). If $g$ is of order a power of $p$, $g$ acts without fixed points on $\bar{X}$ and
\[
\Tr(g^*, H_c^*(X)) = \Tr(g^*, H_c^*(\bar{X}, j_!\overline{\mathbb{Q}}_\ell)) = 0
\]
by the Lefschetz fixed point theorem applied to $\bar{X}$ and the sheaf $j_!\overline{\mathbb{Q}}_\ell$.

\subsection{3.13} \textit{Historical remark.} The method we have followed, using (3.5), was first used by Zarelua (\textit{On finite groups of transformations}, Proc.\ Int.\ Symp.\ on Topology and its Applications, 1968, pp.\ 334--339) and independently by J.\ L.\ Verdier to prove that if a finite group $G$ acts freely on a topological space $X$ of finite cohomological dimension, and if the $H_c^i(X, \mathbb{Z})$ are of finite type, then $\sum (-1)^i H^i(X, \mathbb{Z}) \otimes \mathbb{Q}$ is a multiple of the regular representation of $G$.

\section{The character formula}

The assumptions (0.1.1) are in force in this chapter and the next four.

\begin{definition}[4.1]
Let $T$ be an $F$-stable maximal torus of $G$. The Green function $Q_{T,G}(u)$ is the restriction to the unipotent elements of the character of the virtual representation $R_T^\theta \subset B$, where $B$ is any Borel subgroup containing $T$.
\end{definition}

This Green function does not depend on $B$ (by (1.6) and (1.17 (i))); it only depends on the $G^F$-conjugacy classes of $T$ and $u$. The natural map $x \mapsto \bar{x}$ from $G$ to its adjoint group induces a bijection on the unipotent sets, and
\begin{equation}\label{eq:4.1.1}
Q_{T,G}(u) = Q_{\bar{T}, \bar{G}}(\bar{u}).
\end{equation}
The Green function is integer-valued (3.3) and is a restriction of a character of $G^F$, hence $Q_{T,G}(u^n) = Q_{T,G}(u)$ if $(n, p) = 1$.

Let $a$ be the smallest integer $\ge 1$ such that $F^{a!}$ is the identity on $W$. From the proof of (3.3) we see that $\sum_i \Tr(F^{a!n}, H_c^i(X_T \subset B))t^n$ is a rational function of $t$ and that
\begin{equation}\label{eq:4.1.2}
Q_{T,G}(u) = \sum_i (-1)^i \Tr(F^{a!n}, H_c^i(X_T \subset B)_u^{\ad u})
\end{equation}

In this chapter, we will express the character of $R_T^\theta \subset B$ in terms of $\theta$ and of Green functions.

\begin{theorem}[4.2]
Let $x = su$ be the Jordan decomposition of $x \in G^F$. Then
\[
\Tr(x, R_T^\theta \subset B) = \frac{1}{|Z^0(s)^F|} \sum_{\substack{g \in G^F \\ \ad g(T) \subset Z^0(s)}} Q_{\ad g T, Z^0(s)}(u) \cdot \ad g(\theta)(s).
\]
\end{theorem}

Let $\varepsilon_x$ be the function on $G^F$ whose value is 1 at $x \in G^F$ and 0 elsewhere. One can rewrite 4.2 as giving the following formula for the character of $R_T^\theta \subset B$:
\begin{equation}\label{eq:4.2.1}
\Tr(\ ,\ R_T^\theta \subset B) = \sum_{t \in T^F} \frac{1}{|Z^0(t)^F|} \theta(t) \sum_{\substack{u \in Z^0(t)^F \\ \text{unipotent}}} Q_{T, Z^0(t)}(u) \varepsilon_{\ad g(tu)}
\end{equation}
where we put $Q_{T, Z^0(t)}(u) = 0$ whenever $u$ is not unipotent.

\begin{corollary}[4.3]
$R_T^\theta \subset B$ is independent of the choice of $B$ ($B \supset T$).
\end{corollary}

From now on we shall write $R_T^\theta$ for $R_T^\theta \subset B$. We will deduce 4.2 from 3.10 and the following geometrical facts.

\begin{proposition}[4.4]
\begin{enumerate}
\item[(i)] If a Borel subgroup $B_1$ of $G$ contains $s$, then $B_1 \cap Z^0(s)$ is a Borel subgroup of $Z^0(s)$.
\item[(ii)] Pick $T \subset B$ in $G$. Any Borel subgroup $B_1$ such that $s \in B_1$ is of the form $\ad g B$ with $\ad g T \subset Z^0(s)$ (i.e., $g^{-1}sg \in T$). The left coset $\tau_{T,B}(B_1) \overset{\text{def}}{=} Z^0(s) \cdot g$ depends only on $T$, $B$, and $B_1$.
\item[(iii)] The map $B_1 \mapsto B_1 \cap Z^0(s)$ is an isomorphism from the space of Borel subgroups containing $s$, such that $\tau_{T,B}(B_1) = \tau_{T,B}(B_1')$ and the space of Borel subgroups of $Z^0(s)$.
\end{enumerate}
\end{proposition}

\textit{Proof.} (i) is well known. Put $B_1 = \ad g_1 B$. We have $g_1^{-1}sg_1 \in B$ hence there exists $u$ in the unipotent radical of $B$ such that $u^{-1}g_1^{-1}sg_1u \in T$. We take $g = g_1u$. Put $T' = \ad g T$. If $g' = hg$ is also such that $\ad g' B = B_1$ (resp.\ $\ad g' T' \subset Z^0(s)$) then $h \in B_1$ (resp.\ $h \in Z^0(s) \cdot N(T')$, by the conjugacy of maximal tori in $Z^0(s)$). If $W = N(T')/T'$ is the Weyl group of $G$ and $W_1 = (N(T') \cap Z^0(s))/T'$ that of $Z^0(s)$, the Bruhat decomposition for $Z^0(s)$ reads
\[
Z^0(s) = (B_1 \cap Z^0(s)) \cdot W_1 \cdot (B_1 \cap Z^0(s)),
\]
and
\[
(B_1 \cap Z^0(s)) \cdot N(T') = B_1 \cdot W_1 \cdot (B_1 \cap Z^0(s)) \cap N(T') \subset B_1 \cdot W_1^0 \cdot B_1 \cap N(T') = W_1^0 \cdot T' \subset Z^0(s).
\]
Hence
\[
B_1 \cap (Z^0(s) \cdot N(T')) \subset Z^0(s)
\]
and
\[
B_1 \cap (Z^0(s) \cdot N(T')) = B_1 \cap Z^0(s).
\]
The element $g$ such that $\ad g B = B_1$ and $\ad g T \subset Z^0(s)$ is unique modulo $B_1 \cap Z^0(s)$, hence (ii) and (iii).

A $Z^0(s)$-left coset $z = Z^0(s) \cdot g$, with $\ad g T \subset Z^0(s)$, defines an isomorphism $\gamma$ from the maximal torus of $G$ to that of $Z^0(s)$; $\ad g$ maps $T$ (contained in the Borel subgroup $B$) to $\ad g T \subset Z^0(s)$ (contained in the Borel subgroup $\ad g B \cap Z^0(s)$). This construction will allow us to compare the relative position of Borel subgroups in $G$ and $Z^0(s)$.

\begin{proposition}[4.5]
Let $B_1'$ and $B_1''$ be Borel subgroups containing $s$, with $\tau_{T,B}(B_1') = \tau_{T,B}(B_1'')$; let $w$ be the relative position of $B_1'$ and $B_1''$ (in $G$) and $w_1$ that of $B_1' \cap Z^0(s)$ and $B_1'' \cap Z^0(s)$ (in $Z^0(s)$). Then $w_1 = w$.
\end{proposition}

Replacing $(T, B)$ by a conjugate, we may assume that $B = B_1'$ and that $T \subset B_1' \cap Z^0(s) \cap B_1''$. We will identify the maximal torus of $G$ and that of $Z^0(s)$ with $T$, using $T \subset B$ and $T \subset (B \cap Z^0(s))$. We then have $\tau = \id$. Replacing $B_1''$ by a $N(T) \cap Z^0(s)$-conjugate, we may further assume that $w_1 = 1$, i.e., $B_1' \cap Z^0(s) = B_1'' \cap Z^0(s)$. We must prove that $w = 1$, which is clear.

\begin{lemma}[4.6]
If $B$ and $B'$, containing $T$, are in the same relative position as $B_1$ and $B_1'$ containing $s$, then $\tau_{T,B}(B_1) = \tau_{T,B}(B_1')$.
\end{lemma}

Replacing $(B, T, B')$ by a conjugate, we may again assume that $B = B_1'$ and that $T \subset B_1 \cap Z^0(s) \cap B_1''$. In this case $B_1' = B_1''$.

To prove (4.2) we first compute the space $X_T \subset B$ of Borel subgroups $B_1$ containing $s$ and such that $B_1$ and $FB_1$ are in the same relative position as $B$ and $FB$.

\begin{proposition}[4.7]
Let $X_T \subset B(g)$ be the subspace of $X_T \subset B$ consisting of those $B_1$ with $\tau_{T,B}(B_1) = Z^0(s) \cdot g$. Then
\begin{equation}\label{eq:4.7.1}
X_T \subset B = \coprod_{\substack{g \in Z^0(s)^F\backslash G^F \\ \ad g(T) \subset Z^0(s)}} X_T \subset B(g)
\end{equation}
and, for $g \in G^F$, the map $B_1 \mapsto B_1 \cap Z^0(s)$ induces an isomorphism
\begin{equation}\label{eq:4.7.2}
X_T \subset B(g) \xrightarrow{\sim} X_{\ad g T \subset \ad g B \cap Z^0(s)}.
\end{equation}
\end{proposition}

By 4.6, we have $F\tau_{T,B}(B_1) = \tau_{T, FB}(FB_1) = \tau_{T,B}(B_1)$. By Lang's Theorem, the $Z^0(s)$-principal homogeneous space $\tau_{T,B}(B_1)$, being $F$-stable, has a rational point, hence \eqref{eq:4.7.1}. By (4.4 (iii)) and (4.5), the map \eqref{eq:4.7.2} is an isomorphism from $X_T \subset B(g)$ to some $X(w)$ (relative to $Z^0(s)$); to check which $w$ appears, we observe that $\ad g B \in X_T \subset B(g)$.

\textit{Proof of 4.2.} The theorem is now an immediate application of (4.7) and (3.10) applied to the action (1.17) of $G^F \times T^F$ on $\dot{X}_T \subset B$.

\section{Characters of tori}

We will assume chosen isomorphisms
\begin{equation}\label{eq:5.0.1}
k^* \xrightarrow{\sim} (\mathbb{Q}/\mathbb{Z})_{p'}
\end{equation}
and
\begin{equation}\label{eq:5.0.2}
(\text{roots of unity of order prime to }p\text{ in }\overline{\mathbb{Q}}_\ell^*) \xrightarrow{\sim} (\mathbb{Q}/\mathbb{Z})_{p'}.
\end{equation}

\subsection{5.1} Let $T$ be a torus over $k$. Besides the character group $X(T) = \Hom(T, \mathbb{G}_m)$, we will consider its dual $Y(T) = \Hom(\mathbb{G}_m, T)$. The duality is given by
\[
\langle x, h \rangle = \text{the degree of }x \circ h \in \Hom(\mathbb{G}_m, \mathbb{G}_m) = \mathbb{Z}
\]
($x \in X(T)$ and $h \in Y(T)$). The maps $(h, x) \mapsto h(x)$ and $t \mapsto (x \mapsto x(t))$ induce isomorphisms
\begin{equation}\label{eq:5.1.1}
Y(T) \otimes k^* = T(k) = \Hom(X(T), k^*).
\end{equation}

If $T$ is a maximal torus in $G$, the roots of $T$ in $G$ are elements of $X(T)$, while the coroots belong to $Y(T)$. If $\alpha$ is a root, the corresponding coroot $H_\alpha$ is characterised as follows: there is a homomorphism $u: SL(2) \to G$, which maps the subgroup $\begin{pmatrix} 1 & * \\ 0 & 1 \end{pmatrix}$ onto the root subgroup $U_\alpha$, and whose restriction to the group of diagonal matrices (identified with $\mathbb{G}_m$ by $x \mapsto \begin{pmatrix} x & 0 \\ 0 & x^{-1} \end{pmatrix}$) is $H_\alpha$.

\subsection{5.2} Using the chosen isomorphism \eqref{eq:5.0.1}, we can rewrite \eqref{eq:5.1.1} as
\begin{equation}\label{eq:5.2.1}
T(k) = Y(T) \otimes (\mathbb{Q}/\mathbb{Z})_{p'} = (Y(T) \otimes \mathbb{Q}/Y(T))_{p'}.
\end{equation}
If $T$ is obtained by extension of scalars from a torus $T_0/\mathbb{F}_q$, the Frobenius map $F$ induces a map $F: Y(T) \to Y(T)$ ($Y$ is a covariant functor) and $T^F$ is the subgroup $\Ker(F - 1)$ of $T(k) = Y(T) \otimes (\mathbb{Q}/\mathbb{Z})_{p'}$; since $F - 1$ is divisible by $p$, it is also the kernel of $F - 1$ in $Y(T) \otimes \mathbb{Q}/\mathbb{Z}$:
\begin{equation}\label{eq:5.2.2}
0 \longrightarrow T^F \longrightarrow Y(T) \otimes \mathbb{Q}/\mathbb{Z} \xrightarrow{F-1} Y(T) \otimes \mathbb{Q}/\mathbb{Z} \longrightarrow 0.
\end{equation}
By applying the Snake Lemma to
\[
\begin{array}{ccccccccc}
0 & \longrightarrow & Y(T) & \longrightarrow & Y(T) \otimes \mathbb{Q} & \longrightarrow & Y(T) \otimes \mathbb{Q}/\mathbb{Z} & \longrightarrow & 0 \\
& & \downarrow{F-1} & & \downarrow{F-1} & & \downarrow{F-1} & & \\
0 & \longrightarrow & Y(T) & \longrightarrow & Y(T) \otimes \mathbb{Q} & \longrightarrow & Y(T) \otimes \mathbb{Q}/\mathbb{Z} & \longrightarrow & 0
\end{array}
\]
we get another exact sequence
\begin{equation}\label{eq:5.2.3}
0 \longrightarrow Y(T)_{F-1} \longrightarrow Y(T) \longrightarrow T^F \longrightarrow 0.
\end{equation}
The maps in (5.2.2), (5.2.3) depend on the choice \eqref{eq:5.0.1}. An intrinsic expression for these sequences would be as follows:

$T^F$ is the image of the middle map in the exact sequence
\[
0 \longrightarrow Y(T) \otimes \mathbb{Z}_{p'}(1) \xrightarrow{F-1} Y(T) \otimes \mathbb{Z}_{p'}(1) \longrightarrow Y(T) \otimes \mathbb{Z}_{p'}(1) \otimes \mathbb{Q}/\mathbb{Z} \longrightarrow 0.
\]

Let us map the sequences (5.2.2), (5.2.3) into $\mathbb{Q}/\mathbb{Z}$. The isomorphism \eqref{eq:5.0.2} provides an isomorphism
\[
(T^F)^\wedge \overset{\text{(def)}}{=} \Hom(T^F, \overline{\mathbb{Q}}_\ell^*) \xrightarrow{\sim} \Hom(T^F, \mathbb{Q}/\mathbb{Z}),
\]
and we get exact sequences
\begin{equation}\label{eq:5.2.2*}
0 \longrightarrow X(T) \xrightarrow{F-1} X(T) \longrightarrow (T^F)^\wedge \longrightarrow 0
\end{equation}
\begin{equation}\label{eq:5.2.3*}
0 \longrightarrow (T^F)^\wedge \longrightarrow X(T) \otimes \mathbb{Q}/\mathbb{Z} \xrightarrow{F-1} X(T) \otimes \mathbb{Q}/\mathbb{Z} \longrightarrow 0.
\end{equation}
The dual torus $T^*$ of $T$ is defined by the rule $X(T^*) = Y(T)$, (hence $Y(T^*) = X(T)$); its $\mathbb{F}_q$-structure is defined by ($F$ on $Y(T^*)$) = (${}^tF$ on $X(T)$). Via some isomorphism
\begin{equation}\label{eq:5.2.4}
(T^F)^\wedge \xrightarrow{\sim} T^{*F}
\end{equation}
the sequences (5.2.3), (5.2.2) for $T^*$ are identical with \eqref{eq:5.2.2*}, \eqref{eq:5.2.3*}.

\subsection{5.3} If we go from $\mathbb{F}_q$ to $\mathbb{F}_{q^n}$, $F$ is replaced by $F^n$. Composition with the norm map
\[
N: T^{F^n} \longrightarrow T^F, \quad t \mapsto t \cdot F(t) \cdots F^{n-1}(t)
\]
is an injection $^tN: (T^F)^\wedge \to (T^{F^n})^\wedge$. We have commutative diagrams
\[
\begin{array}{ccccccccccc}
& & & & F^n-1 & & & & & & \\
0 & \longrightarrow & Y(T) & \longrightarrow & Y(T) & \longrightarrow & T^{F^n} & \longrightarrow & Y(T) \otimes \mathbb{Q}/\mathbb{Z} & \longrightarrow & 0 \\
& & \downarrow{N} & & \downarrow{N} & & \downarrow{N = \prod_0^{n-1}F^i} & & \downarrow{N} & & \\
0 & \longrightarrow & X(T) & \longrightarrow & X(T) & \longrightarrow & (T^F)^\wedge & \longrightarrow & X(T) \otimes \mathbb{Q}/\mathbb{Z} & \longrightarrow & 0 \\
& & & & F-1 & & & & & & \\
& & & & \downarrow & & \downarrow & & & & \\
& & & & X(T) & \longrightarrow & (T^F)^\wedge & \longrightarrow & X(T) \otimes \mathbb{Q}/\mathbb{Z} & \longrightarrow & 0
\end{array}
\]
where $X$ and $Y$ stand for $X(T)$ and $Y(T)$. In particular, for $\theta$ a character of $T^F$, $\theta$ and $\theta \circ N$ have the same image in $X(T) \otimes \mathbb{Q}/\mathbb{Z}$, and (5.2.4) gives rise to a commutative diagram
\[
\begin{array}{ccc}
(T^F)^\wedge & \xrightarrow{\sim} & T^{*F} \\
\downarrow{^tN} & & \downarrow{N} \\
(T^{F^n})^\wedge & \xrightarrow{\sim} & T^{*F^n}
\end{array}
\]

\begin{proposition}[5.4]
Let $T$ and $T'$ be two $F$-stable maximal tori in $G$, and let $\theta$, $\theta'$ be characters of $T^F$ and $T'^F$. We identify them with characters of $Y(T)$ and $Y(T')$, by (5.2.3). The following conditions are equivalent:
\begin{enumerate}
\item[(i)] For some $x \in G$, with $\ad x(T) = T'$, the map induced by $\ad x: Y(T) \to Y(T')$ carries $\theta$ to $\theta'$;
\item[(ii)] For some $n$, the pairs $(T, \theta \circ N)$, $(T', \theta' \circ N)$, where $N$ is the norm from $T^{F^n}$ to $T^F$ (resp.\ $T'^{F^n}$ to $T'^F$) are $G^{F^n}$-conjugate.
\end{enumerate}
\end{proposition}

By 5.3, the condition (i) is invariant under the replacement of $F$ by $F^n$ and of $\theta$, $\theta'$ by $\theta \circ N$, $\theta' \circ N$. It hence suffices to check (5.4) in the trivial case where $T$ and $T'$ are split.

\begin{definition}[5.5]
The pairs $(T, \theta)$, $(T', \theta')$ are said to be \textit{geometrically conjugate} when the equivalent conditions of (5.4) hold.
\end{definition}

\subsection{5.6} Let $(T, \theta)$ be as above. The choice of a Borel subgroup $B$ containing $T$ defines an isomorphism of $T$ with the torus $T$; the corresponding isomorphism $Y(T) \xrightarrow{\sim} Y(T)$ carries $\theta$ to a character of $Y(T)$ with values in the roots of unity of order prime to $p$ in $\overline{\mathbb{Q}}_\ell$. The isomorphism \eqref{eq:5.0.2} identifies it with an element of $X(T) \otimes (\mathbb{Q}/\mathbb{Z})_{p'}$. Its class $[\theta]$ in $(X(T) \otimes (\mathbb{Q}/\mathbb{Z})_{p'})^W$ is independent of the choice of $B$. It is invariant under $F$. Let us put
\[
\mathcal{S} = [(X(T) \otimes (\mathbb{Q}/\mathbb{Z})_{p'})/W]^F \xrightarrow{\sim} [(X(T) \otimes \mathbb{Q}/\mathbb{Z})/W]^F.
\]

\begin{proposition}[5.7]
\begin{enumerate}
\item[(i)] The map $\theta \mapsto [\theta]$ induces a bijection from the set of geometric conjugacy classes of pairs $(T, \theta)$ to $\mathcal{S}$.
\item[(ii)] (Compare Steinberg \cite[ p.\ 93]{Stein}). The number of geometric conjugacy classes of pairs $(T, \theta)$ is $|Z^{0F}| \cdot q^l$ where $Z^0$ is the identity component of the centre of $G$ and $l$ is the semisimple rank of $G$.
\end{enumerate}
\end{proposition}

(i) The injectivity is clear on 5.4(i). Let us prove surjectivity. If the class mod $W$ of $x \in X(T) \otimes (\mathbb{Q}/\mathbb{Z})_{p'}$ is invariant by $F$, one indeed has $t(wF)x = x$ for some $w \in W$, and $x$ corresponds to a character of $T(w)^F$, hence of $T^F$, for some $F$-stable maximal torus $T$ in $G$ (1.14).

(ii) As in (5.2), we may now identify geometric conjugacy classes of pairs $(T, \theta)$ with $\mathbb{F}_q$-rational points of the quotient $T^*/W$. Our task is to compute
\[
|(T^*/W)^F| = \sum_i (-1)^i \Tr(F^*, H_c^i(T^*/W)) = \sum_i (-1)^i \Tr(F^*, H_c^i(T^*)^W).
\]
Let $G'$ be the derived group of $G$; the torus $T^*$ is isogenous to (hence has the same cohomology as) the product of $Z^{0*}$ with $T'^*$ ($T'$ being the torus of $G'$). Since any torus has as many rational points as its dual, we get by K\"unneth's Theorem
\[
|(T^*/W)^F| = |Z^{0F}| \sum_i (-1)^i \Tr(F^*, H_c^i(T'^*)^W).
\]
We have
\[
H^*(T'^*, \overline{\mathbb{Q}}_\ell) = \Lambda^* H^1(T'^*, \overline{\mathbb{Q}}_\ell) = \Lambda^* (Y(T') \otimes \overline{\mathbb{Q}}_\ell(-1))
\]
hence by Poincar\'e duality,
\[
H_c^{2l-i}(T'^*, \overline{\mathbb{Q}}_\ell) \cong \Hom(\Lambda^i Y(T'), \overline{\mathbb{Q}}_\ell(i-l))
\]
and
\[
\Tr(F^*, H_c^{2l-i}(T'^*)^W) = q^{l-i} \Tr(F, \Lambda^i Y(T')^W).
\]
The next lemma shows that this trace is zero for $i = 0$, hence the $q^l$ factor.

\begin{lemma}[5.8]
Let $W \subset \Aut(V)$ be a finite group generated by reflections of real vector space $V$ of dimension $d$. We assume that $V^W = 0$. Then $(\Lambda^i V)^W = 0$ for $i \neq 0$.
\end{lemma}

For a proof, see Bourbaki \cite[Ch.\ V, Ex.\ 3 of \S 2]{Bou}.

\subsection{5.9} Let $T$ be an $F$-stable maximal torus in $G$ and let $\alpha$ be a root of $T$. For $\theta$ a character of $T^F$, we will say that $\theta$ is \textit{orthogonal} to $H_\alpha$: $\langle H_\alpha, \theta \rangle = 0$ if $\theta$, viewed as a character of $Y(T)$ (5.2.3), is identically $1$ on $H_\alpha$. The character $\theta$ of $Y(T)$ is then invariant by the corresponding reflection $s_\alpha$.

\subsection{5.10} Let $\pi: \tilde{G} \to G$ be the simply connected covering of the derived group of $G$, and let $\tilde{T}$ be the inverse image of $T$ in $\tilde{G}$. The group $\tilde{T}$ is connected: it is a maximal torus in $\tilde{G}$. It acts on $\tilde{T} \times G$ by $\tilde{T} \cdot (t, g) = (tw(t)^{-1}, tg)$, and the map $(t, g) \mapsto tg$ induces an isomorphism $\tilde{T}\backslash \tilde{T} \times G \xrightarrow{\sim} G$, hence (Lang's Theorem) $\tilde{T}^F\backslash \tilde{T}^F \times G^F \xrightarrow{\sim} G^F$. We have
\begin{equation}\label{eq:5.10.1}
T^F/\pi(\tilde{T}^F) \xrightarrow{\sim} G^F/\pi(\tilde{G}^F).
\end{equation}

\begin{proposition}[5.11]
\begin{enumerate}
\item[(i)] A character of $T^F$ is the restriction to $T^F$ of a character of $G^F/\pi(\tilde{G}^F)$ if and only if it is orthogonal to all coroots.
\item[(ii)] Let $\theta$ be a character of $G^F/\pi(\tilde{G}^F)$, $T$ and $T'$ two $F$-stable maximal tori and let $x \in G$ be such that $\ad x(T) = T'$. Then $\ad x: Y(T) \to Y(T')$ carries the restriction of $\theta$ to $T^F$ (viewed as a character of $Y(T)$) onto the restriction of $\theta$ to $T'^F$: the restrictions of $\theta$ to the $F$-stable maximal tori are all geometrically conjugate.
\end{enumerate}
\end{proposition}

The assertion (i) follows from \eqref{eq:5.10.1}, and the fact that $Y(\tilde{T}) \subset Y(T)$ is spanned by the coroots $H_\alpha$. To prove (ii), we will use the criterion 5.4 (ii). Let $n$ be such that $x \in G^{F^n}$. We will prove that if $y \in T^{F^n}$, $N(y)$ and $N(\ad x(y))$ have the same image in $G^F/\pi(\tilde{G}^F)$, hence that $\ad x$ transforms $(\theta \mid T^F) \circ N$ into $(\theta \mid T'^F) \circ N$.

The map $t \mapsto N(t)^{-1}N(\ad x(t)): T^{F^n} \to G$ lifts uniquely into a map $\varphi: T^{F^n} \to \tilde{G}$ mapping $e$ to $e$. Indeed,

(a) it factors through $T/Z$;

(b) the map $t \mapsto N(t)^{-1}N(\ad x(t)): T^{F^n} \to G$ factors through $T/Z = \tilde{T}/\tilde{Z}$ and the desired lifting is the composite $T^{F^n} \to T/Z = \tilde{T}/\tilde{Z} \to \tilde{G}$.

Similarly, the map
\[
(a, t) \mapsto \theta(a)^{-1}N(t)^{-1}N(\ad x(t))\ad x(a): T^{F^n} \times T^{F^n} \longrightarrow G
\]
lifts into $\Phi: T^{F^n} \times T^{F^n} \to \tilde{G}$, with $\Phi(e, t) = \varphi(t)$. The identity
\begin{align*}
F(N(t)^{-1}N(\ad x(t))) &= N(Ft)^{-1}N(F\ad x(t)) \\
&= (t^{-1}F^nt)^{-1}N(t)^{-1}N(\ad x(t))\ad x(t^{-1}Ft) \\
&= (t^{-1}F^nt)^{-1}N(t)^{-1}N(\ad x(t))\ad x(t^{-1}Ft)
\end{align*}
lifts into
\[
F\varphi(t) = \Phi(t^{-1}F^nt, t)^{-1}\varphi(t)\Phi(t^{-1}F^nt, t).
\]
Putting $t = y$, we have $y^{-1}F^ny = e$, hence $F\varphi(y) = \varphi(y)$ and $\varphi(y) \in \tilde{G}^F$. We have $N(y)^{-1}N(\ad x(y)) \in \pi(\tilde{G}^F)$, as required.

\begin{remark}[5.12]
The argument above can be used to show the existence of a norm map $N: G^{F^n}/\pi(\tilde{G}^{F^n}) \to G^F/\pi(\tilde{G}^F)$ which for each $F$-stable maximal torus $T$ gives rise to a commutative diagram
\[
\begin{array}{ccc}
T^{F^n}/\pi(\tilde{T}^{F^n}) & \longrightarrow & G^{F^n}/\pi(\tilde{G}^{F^n}) \\
\downarrow{N} & & \downarrow{N} \\
T^F/\pi(\tilde{T}^F) & \longrightarrow & G^F/\pi(\tilde{G}^F)
\end{array}
\]
\end{remark}

\begin{theorem}[5.13]
Let $(T, \theta)$ be as above. We use \eqref{eq:5.2.3} to identify $\theta$ with a character of $Y(T)$. If the centre of $G$ is connected, then the stabilizer of $\theta$ in the Weyl group $W = N(T)/T$ is generated by the reflections $s_{H_\alpha}$ for $H_\alpha$ a coroot orthogonal to $\theta$.
\end{theorem}

The key point in the proof is the same as in Steinberg's Theorem that if the derived group $\tilde{G}$ is simply connected, then the centralizer of any semi-simple element is connected.

Let us identify $\theta$ with the corresponding element in $X(T) \otimes (\mathbb{Q}/\mathbb{Z})_{p'}$ (invariant under $F$). Let $\theta_1 \in X(T) \otimes \mathbb{Q}$ be a representative for it. This $\theta_1$ has no $p$ in the denominator. We have the dictionary:
\begin{enumerate}
\item[(a)] $\langle H_\alpha, \theta \rangle = 0 \iff \langle H_\alpha, \theta_1 \rangle \in \mathbb{Z}$;
\item[(b)] $\theta$ is fixed by $w \in W \iff$ for some $x \in X(T)$, $w\theta_1 + x = \theta_1$.
\end{enumerate}

Let $X_{\ad}$ be the subgroup of $X$ generated by the roots. It is the character group of the image $T_{\ad}$ of $T$ in the adjoint group. The character group of the centre $Z = \Ker(T \to T_{\ad})$ of $G$ is $X/X_{\ad}$; the character group of $Z/Z_{\red}$ is the torsion subgroup of $X/X_{\ad}$, hence
\begin{enumerate}
\item[(c)] $Z$ is connected and smooth (resp.\ connected) if and only if $X/X_{\ad}$ is torsion free (resp.\ has no torsion prime to $p$).
\end{enumerate}

For each root $\alpha$, $s_\alpha(\theta_1) = \theta_1 - \langle H_\alpha, \theta_1 \rangle \alpha$. The number $\langle H_\alpha, \theta_1 \rangle$ has no $p$ in the denominator; hence $\langle H_\alpha, \theta_1 \rangle \cdot x \in X(T)$ if and only if $\langle H_\alpha, \theta_1 \rangle \in \mathbb{Z}$; $s_\alpha(\theta) = \theta$ if and only if $\langle H_\alpha, \theta \rangle = 0$. It remains to check that the stabilizer of $\theta$ is generated by the reflections it contains.

Fix $N$ such that $(1/p^N)X_{\ad} \supset X \cap (X_{\ad} \otimes \mathbb{Q})$ and such that $p^N \equiv 1 \pmod{\text{the order of }\theta}$. For any $w \in W$, $w\theta_1 - \theta_1$ is in $X_{\ad} \otimes \mathbb{Q}$, hence if $w\theta_1 - \theta_1$ is in $X$, $w(p^N\theta_1) - (p^N\theta_1) \in X_{\ad}$. Replacing $\theta_1$ by $p^N\theta_1$, which is also a representative of $\theta$, we may assume that $\theta$ is fixed by $w \in W$ if and only if $w\theta_1 + x = \theta_1$ for some $x \in X_{\ad}$. The stabilizer of $\theta$ in $W$ is the image in $W$ of the stabilizer of $\theta_1$ in the affine Weyl group. It remains to apply Bourbaki \cite[Ch.\ VI, Ex.\ 1 of \S 2]{Bou}.

\begin{remark}[5.14]
The proof shows that, if we use a suitable Borel subgroup $B \supset T$ to identify $W$ with $W$, the stabilizer of $\theta$ becomes the subgroup of $W$ generated by some of the reflections corresponding either to a simple root or to the negative of the highest coroot.
\end{remark}

\begin{definition}[5.15]
\begin{enumerate}
\item[(i)] The character $\theta$ of $T^F$ is \textit{nonsingular} if it is not orthogonal to any coroot.
\item[(ii)] $\theta$ is \textit{in general position} if it is not kept fixed by any non-trivial element of $(N(T)/T)^F$.
\end{enumerate}
\end{definition}

\begin{proposition}[5.16]
If the centre of $G$ is connected, a character is non-singular if and only if it is in general position.
\end{proposition}

The stabilizer of $\theta$ in $N(T)/T$ is a group generated by reflections, and is stable by Frobenius. We have to prove that the subgroup fixed by Frobenius is non-trivial. This follows from the next lemma.

\begin{lemma}[5.17]
Let $V \neq \{0\}$ be a euclidean vector space, and let $W$ be a finite group generated by orthogonal reflections. We assume that $V^W \neq \{0\}$. If $a$ belongs to the normalizer $A$ of $W$ in $O(V)$, the centralizer $W_a$ of $a$ in $W$ is non-trivial.
\end{lemma}

(a) Reduction to the irreducible case. Let $V = \bigoplus V_i$ be the decomposition of $V$ as a direct sum of irreducible root systems; we have $W = \prod W_i$. The automorphism $a$ permutes the $V_i$; by taking a direct factor, we may assume that it permutes them cyclically: $V = \bigoplus_{\mathbb{Z}/r\mathbb{Z}} V_i$ and $aV_i = V_{i+1}$. An element $w = (w_i) \in W = \prod W_i$ is in $W_a$ if and only if $w_i = \ad a^i(w_0)$ and $w_0$ is fixed by $a^r$; we are reduced to proving that $W_{a^r}$ is non-trivial.

(b) In the irreducible case, one always has either $A = W \cup -W$ or $-1 \in W$; both cases are clear.

\begin{corollary}[5.18]
For any $G$, if $\theta$ is in general position, then $\theta$ is non-singular.
\end{corollary}

Let us embed $G$ in a group $G_1$ with connected centre and the same derived group. For instance, one can take $G_1 = G \times T/\{(z, z^{-1}) \mid z \in Z(G)\}$. The torus $T$ is then contained in an $F$-stable maximal torus $T_1$ of $G_1$, and $\theta$ is the restriction to $T^F$ of some character $\theta_1$ of $T_1^F$. If $\theta$ is in general position, $\theta_1$ is so a fortiori, hence non-singular (5.16). A character $\theta_1$ of $T_1^F$ is non-singular if and only if its restriction to $T^F$ is, hence (5.18).

When all roots of $G$ have the same length, geometric conjugacy can be given the following convenient description.

\begin{definition}[5.19]
(Assuming all roots to be of the same length) Let $T$ be an $F$-stable maximal torus and let $\theta$ be a character of $T^F$. The connected centralizer $S(T, \theta)$ of $(T, \theta)$ in $G$ is the $F$-stable reductive subgroup of $G$ with maximal torus $T$ and with root subgroups relative to $T$ the $U_\alpha$ for which $\langle H_\alpha, \theta \rangle = 0$.
\end{definition}

If $\langle H_\alpha, \theta \rangle = \langle H_\beta, \theta \rangle = 0$ and $\alpha + \beta$ is a root, then $H_{\alpha+\beta} = H_\alpha + H_\beta$, hence $\langle H_{\alpha+\beta}, \theta \rangle = 0$ and the definition makes sense. Let $\pi: \widetilde{S(T, \theta)} \to S(T, \theta)$ be the simply connected covering of the derived subgroup of $S(T, \theta)$. By 5.11 (i), $\theta$ extends to a character $\theta_S$ of $S(T, \theta)^F/\pi(\widetilde{S(T, \theta)}^F)$.

\begin{proposition}[5.20]
If the centre of $G$ is connected and all roots are of the same length, then a pair $(T', \theta')$ is geometrically conjugate to $(T, \theta)$ if and only if $(S(T, \theta), \theta_S)$ and $(S(T', \theta'), \theta'_{S'})$ are $G^F$-conjugate.
\end{proposition}

The ``if'' part follows from 5.11. Conversely, let $(T', \theta')$ be geometrically conjugate to $(T, \theta)$: for some $x \in G$, $\ad x(T) = T'$ and $\ad x: Y(T) \to Y(T')$ maps $\theta$ to $\theta'$ (via (5.2.3)). If $F': T' \to T'$ is the Frobenius map of $T'$, $\ad x^{-1}(F') = \ad wF$ for some $w$ in the Weyl group of $T$. We have ${}^tF\theta = \theta$, and ${}^t(\ad wF)\theta = \theta$, hence $\theta$ is fixed by $w$, which by part (i) belongs to the Weyl group of $T$ in $S(T, \theta)$. Let $T'' = \ad y(T)$, $y \in S(T, \theta)$ be an $F$-stable torus in $S(T, \theta)$, with Frobenius $F''$, such that $(\ad y)^{-1}(F'') = (\ad x)^{-1}(F')$. Replacing $(T, \theta)$ by $(T'', \ad y(\theta))$ (by applying 5.11) and replacing $x$ by $xy^{-1}$, we are reduced to the case where $\ad x F = F' \ad x$. In this case, $(T, \theta)$ and $(T', \theta')$ are $G^F$-conjugate; the identity $\ad x(Ft) = F\ad x(t)$ for $t \in T$ amounts to $x^{-1}Fx \in T$; and replacing $x$ by $xt^{-1}$ with $t \in T$, $t^{-1}Ft = x^{-1}Fx$, makes $x$ rational.

When roots are not all of the same length, one has to work in the dual group $G^*$ of $G$.

\begin{definition}[5.21]
A group $G^*$ dual to $G$ is a reductive group $G^*$ defined over $\mathbb{F}_q$, whose maximal torus $T^*$ is provided with an isomorphism with the dual of the maximal torus $T$ of $G$, this isomorphism carrying simple roots to simple coroots.
\end{definition}

Here are some properties of this duality. Let $G$ and $G^*$ be dual.

\textbf{(5.21.1)} $G$ and $G^*$ have the same Weyl group $W$. The action of $W$ on $Y(T^*)$ is the contragredient of its action on $Y(T)$.

\textbf{(5.21.2)} However, Frobenius does not act in the same way on $W$ in $G$ and $G^*$: it acts in inverse ways (the origin of this is that $F$ on $Y(T^*)$ is the transpose of $F$ on $Y(T)$).

\textbf{(5.21.3)} The conjugacy classes of pairs $(T, B)$ ($T$ an $F$-stable maximal torus, $B$ a Borel subgroup containing it) in $G$ and $G^*$ correspond. With the notations of (1.13), (1.14), to the class of $(T, B)$ we associate the class of $(T^*, B^*)$ with ${}^t(\ad h(T, B)^{-1} \circ F) = \ad h(T^*, B^*)^{-1} \circ {}^tF$; the tori $T$ and $T^*$ are in duality.

\textbf{(5.21.4)} This induces a bijection between the rational conjugacy classes of maximal tori in $G$ and $G^*$. If $T$ and $T'$ are in corresponding classes, we have a natural class (mod $(N(T)/T)^F$) of isomorphisms between $T^*$ and $T'$.

\textbf{(5.21.5)} Let $T$ be an $F$-stable torus in $G$, and let $\theta$ be a character of $T^F$. Fix $T'$, a corresponding torus in $G^*$. The character $\theta$ defines an element of $T^{*F}$, then a $(N(T')/T')^F$-conjugacy class of elements $\theta'$ of $T'$. In that way, we get a bijection between $G^F$-conjugacy classes of pairs $(T, \theta)$ as above, and $G^{*F}$-conjugacy classes of pairs $(T', \theta')$, $T'$ an $F$-stable maximal torus of $G^*$ and $\theta'$ an element of $T'^F$. This correspondence is compatible with extension of scalars from $\mathbb{F}_q$ to $\mathbb{F}_{q^n}$ (cf.\ 5.3.).

\textbf{(5.21.6)} By forgetting $T'$, we see that each $G^F$-conjugacy class of pairs $(T, \theta)$ defines an element $\theta' \in G^{*F}$, well-defined up to $G^{*F}$-conjugacy.

\begin{proposition}[5.22]
Two pairs $(T_1, \theta_1)$ and $(T_2, \theta_2)$ are geometrically conjugate if and only if $\theta_1'$ and $\theta_2'$ are geometrically conjugate.
\end{proposition}

An extension of scalars (cf.\ 5.4, 5.5) reduces us to the case where $T_1$ and $T_2$ are split. Conjugating, we may assume further that $T_1 = T_2$. In this case, geometric conjugacy is $W$-conjugacy, and the proposition is clear.

\begin{proposition}[5.23]
Let $G$ and $G^*$ be dual. Then, the centre of $G$ is connected and smooth (resp.\ connected) if and only if the derived group of $G^*$ is simply connected (resp.\ if its simply connected covering is unseparable).
\end{proposition}

Put $Y = Y(T^*)$, $T^*$ the maximal torus of $G^*$; let $Y_1$ be the subgroup generated by the coroots; and put $Y_2 = Y \cap Y_1 \otimes \mathbb{Q}$. Then, $Y_2$ is the $Y$-group of the maximal torus of the derived group, and $Y_1$ that of the maximal torus of its universal covering. The group of this covering is the dual of the Pontrjagin dual of $Y_2/Y_1$, and (5.23) follows by comparison with point (c) in the proof of (5.13).

\begin{corollary}[5.24]
If the centre of $G$ is connected, two pairs $(T_1, \theta_1)$, $(T_2, \theta_2)$ as in (5.22) are geometrically conjugate if and only if $\theta_1'$ and $\theta_2'$ are $G^{*F}$-conjugate.
\end{corollary}

Indeed, by a theorem of Steinberg, geometric conjugacy amounts in this case (5.23) to conjugacy, for semi-simple elements of $G^*$ (their centralizers are connected).

\begin{definition}[5.25]
Let $(T', \theta')$ correspond to $(T, \theta)$ as in (5.21.5). The pair $(T, \theta)$ is \textit{maximally split} if the $\mathbb{F}_q$-rank of $T$ is equal to that of the centralizer of $\theta'$, i.e., if $T'$ is maximally split in that centralizer.
\end{definition}

Any geometric conjugacy class of pairs $(T, \theta)$ contains maximally split pairs.

\begin{proposition}[5.26]
If the centre of $G$ is connected, two geometrically conjugate maximally split pairs $(T_1, \theta_1)$, $(T_2, \theta_2)$ are $G^F$-conjugate.
\end{proposition}

Via (5.21.5) and (5.24), this amounts to the known fact that two maximally split maximal tori in the centralizer $Z(\theta')$ of a semi-simple element $\theta'$ of $G^{*F}$ are conjugate by an element of $Z(\theta')^F$.

\begin{proposition}[5.27]
If $(T, \theta)$ is maximally split, and the image of $T$ in the adjoint group is anisotropic, then $\theta$ is non-singular.
\end{proposition}

Fix $(T', \theta')$ as in (5.21.5). By assumption, $Z^0(\theta')$ contains no non-central $F$-stable split torus, hence $Z^0(\theta') = T'$. The roots of $T'$ in $Z^0(\theta')$ correspond to the coroots of $T$ orthogonal to $\theta$, hence the proposition.

\begin{corollary}[5.28]
If $(T, \theta)$ is maximally split, $T$ is contained in an $F$-stable Levi subgroup $L$ of an $F$-stable parabolic subgroup $P$, and, in $L$, $(T, \theta)$ is non-singular.
\end{corollary}

One chooses $P$ and $L$ so that the image of $T$ in the adjoint group of $L$ is anisotropic. Being maximally split in $G$, $(T, \theta)$ is a fortiori maximally split in $L$, and one applies (5.27).

The following result will be used in the next chapter.

\begin{proposition}[5.29]
Let $T'$ be a subtorus of a torus $T$, with $T$ defined over $\mathbb{F}_q$ (no assumption on $T'$). Let $\theta$ be a character of $T^F$. It is trivial on $T' \cap T^F$ if and only if, when viewed as a character of $Y(T)$, it is trivial on $((F - 1)Y(T') \otimes \mathbb{Q} \cap Y(T))$.
\end{proposition}

The condition is that $\theta(x) = 0$ for
\[
(F - 1)^{-1}(x) \in (Y(T') \otimes \mathbb{Q} + Y(T))
\]
i.e., for
\[
x \in ((F - 1)Y(T') \otimes \mathbb{Q} \cap Y(T)) + ((F - 1)Y(T)).
\]
The character $\theta$ being trivial on $(F - 1)Y(T)$, this means $\theta(x)$ is trivial for $x \in ((F - 1)Y(T') \otimes \mathbb{Q} \cap Y(T))$.

\section{Intertwining numbers}

\subsection{6.1} Let $T$, $T'$ be two $F$-stable maximal tori in $G$ and let $\theta$, $\theta'$ be characters of $T^F$ and $T'^F$. We put $N_G(T, T') = \{g \in G \mid Tg = gT'\}$ and
\[
W_G(T, T') = T\backslash N_G(T, T') = N_G(T, T')/T'.
\]
We will drop the index $G$ if there is no ambiguity. $F$ acts on $W(T, T')$, and
\[
W(T, T')^F = T^F\backslash N(T, T')^F = N(T, T')^F/T'^F.
\]

\begin{theorem}[6.2]
If $\theta^{-1}$ is not geometrically conjugate to $\theta'$, then
\[
\langle H_c^*(X_T \subset B)_\theta \otimes H_c^*(X_{T'} \subset B')_{\theta'}, G^F \rangle = 0,
\]
i.e., if an irreducible representation of $G^F$ occurs in $H_c^*(X_T \subset B)_\theta$, its dual does not occur in $H_c^*(X_{T'} \subset B')_{\theta'}$.
\end{theorem}

The virtual representation $R_{T'}^{\theta'}$ is the dual of $R_T^\theta$ (as can be seen on the character formula 4.2), hence we have as a corollary:

\begin{corollary}[6.3]
If $\theta$ and $\theta'$ are not geometrically conjugate, no irreducible representation of $G^F$ can occur in both virtual representations $R_T^\theta$ and $R_{T'}^{\theta'}$.
\end{corollary}

The proof of 6.2 will make use of the following homotopy argument:

\begin{proposition}[6.4]
Let $H$ be a connected algebraic group acting on a scheme $Y$, separated and of finite type over $k$. For any $h \in H$, the action of $h$ on $H_c^*(Y, \mathbb{Z}/n)$ is trivial.
\end{proposition}

Let $\pr_2$ be the projection of $H \times Y$ onto $H$ and let $f: H \times Y \to H \times Y$ be defined by $f(h, y) = (h, hy)$:
\[
\begin{array}{ccc}
H \times Y & \xrightarrow{f} & H \times Y \\
\pr_2 \downarrow & & \downarrow \pr_2 \\
H & = & H
\end{array}
\]
By the change-of-basis theorem in cohomology with compact support applied to
\[
\begin{array}{ccc}
H \times Y & \longrightarrow & Y \\
\downarrow & & \downarrow \\
H & \longrightarrow & \Spec k,
\end{array}
\]
$R^i\pr_{2*} \mathbb{Z}/n$ is the constant sheaf $H_c^i(Y, \mathbb{Z}/n)$ on $H$. The automorphism $f$ acts on it and, at $h \in H$, it acts the way $h$ acts on $H_c^i(Y, \mathbb{Z}/n)$. At the identity element of $H$, the action of $f$ is trivial. An endomorphism of a constant sheaf over a connected base is constant, hence $f$ acts trivially everywhere, and the proposition is proved.

\begin{corollary}[6.5]
The conclusion of (6.4) holds in $\ell$-adic cohomology.
\end{corollary}

The proof is a passage to limit.

\subsection{6.6} \textit{Proof of 6.2.} On $\dot{X}_T \subset B \times \dot{X}_{T'} \subset B'$ we have commuting actions of $G^F$, $T^F$, and $T'^F$ ($G^F$ acts diagonally). By K\"unneth's Theorem, we have to prove that
\[
H_c^*(\dot{X}_T \subset B \times \dot{X}_{T'} \subset B')_{\theta, \theta'} = 0,
\]
where the symbol $M_{\theta, \theta'}$, applied to any $T^F \times T'^F$-module $M$, means the subspace of $M$, where $T^F$ and $T'^F$ act by $\theta$ and $\theta'$.

Let $U$ (resp.\ $U'$) be the unipotent radical of $B$ (resp.\ $B'$).

Let $S_T \subset B = \{g \in G \mid g^{-1}Fg \in FU\}$, $S_{T'} \subset B' = \{g' \in G \mid g'^{-1}Fg' \in FU'\}$. The unipotent group $(U \cap FU) \times (U' \cap FU')$ acts freely on $S_T \subset B \times S_{T'} \subset B'$ and the orbit space is $\dot{X}_T \subset B \times \dot{X}_{T'} \subset B'$ (see (1.17)). It follows that
\[
H_c^*(\dot{X}_T \subset B \times \dot{X}_{T'} \subset B')(-d) = H_c^*(S_T \subset B \times S_{T'} \subset B')
\]
where $d$ is the dimension of $(U \cap FU) \times (U' \cap FU')$. This isomorphism is compatible with the action of $G^F \times T^F \times T'^F$; moreover this group acts trivially on $\overline{\mathbb{Q}}_\ell(-d)$. Hence it is sufficient to prove that
\[
H_c^*(S_T \subset B \times S_{T'} \subset B')_{\theta, \theta'} = H_c^*(S_T \subset B \times S_{T'} \subset B'/G^F)_{\theta, \theta'} = 0.
\]
The map
\[
(g, g') \longmapsto (x, x', y), \qquad x = g^{-1}Fg, \quad x' = g'^{-1}Fg', \quad y = g^{-1}g'
\]
defines an isomorphism of $S_T \subset B \times S_{T'} \subset B'/G^F$ with
\[
\mathcal{S} = \{(x, x', y) \in FU \times FU' \times G \mid xFy = yx'\}.
\]
Under this isomorphism, $T^F \times T'^F$ acts on $\mathcal{S}$ by the formula
\begin{equation}\label{eq:6.6.1}
(x, x', y) \longmapsto (t^{-1}xt, t'^{-1}x't', t^{-1}yt'), \quad (t, t') \in T^F \times T'^F.
\end{equation}
Let $U'^-$ be opposed to $U'$ with respect to $T'$. Any $g \in G$ can be written uniquely in the form $g = u_g \dot{n}_g u'_g$, with $u_g \in U \cap \dot{n}_g U'^- \dot{n}_g^{-1}$, $\dot{n}_g \in N(T, T')$ (see (6.1)), $u'_g \in U'$; this follows easily from Bruhat's Lemma. For any $w \in W(T, T')$, let $G_w$ be the set of all $g \in G$ such that $\dot{n}_g$ represents $w$. Then $(G_w)_{w \in W(T, T')}$ is a finite partition of $G$ into locally closed subschemes. It has the property (6.6.2) below, expressing the fact that the closure of a Bruhat cell is a union of Bruhat cells.

(6.6.2) For a suitable ordering of $W(T, T')$, the unions $\Sigma_W = \bigcup_{w' \le w} G_{w'}$ are closed.

Let $\mathcal{S}_w = \{(x, x', y) \in \mathcal{S} \mid y \in G_w\}$. Then $(\mathcal{S}_w)_{w \in W(T, T')}$ is a finite partition of $\mathcal{S}$ into locally closed subschemes, stable under $T^F \times T'^F$. It inherits a property analogous to (6.6.2), that is, the unions $\bigcup_{w' \le w} \mathcal{S}_{w'}$ are closed for any $w$.

The spectral sequence associated to the filtration of $\mathcal{S}$ by these unions shows that, in order to prove that $H_c^*(\mathcal{S})_{\theta, \theta'} = 0$, it is sufficient to prove that $H_c^*(\mathcal{S}_w)_{\theta, \theta'} = 0$ for any $w \in W(T, T')$.

Let
\[
H_w = \{(t, t') \in T \times T' \mid t'F(t')^{-1} = F(\dot{w})^{-1}tF(t)^{-1}F(\dot{w})\}
\]
where $\dot{w} \in N(T, T')$ represents $w$. Then $H_w$ is a closed subgroup of $T \times T'$, containing $T^F \times T'^F$. We define an action of $H_w$ on $\mathcal{S}_w$ as follows. Let $(t, t') \in H_w$; for $(x, x', y) \in \mathcal{S}_w$ define
\[
f_{t,t'}(x, x', y) = (\tilde{x}, \tilde{x}', \tilde{y})
\]
where
\begin{align*}
\tilde{x} &= t^{-1}xF(u_y)tF(t^{-1}u_y^{-1}t), \\
\tilde{x}' &= t'^{-1}x'F(u'_y)^{-1}t'F(t'^{-1}u'_yt'), \\
\tilde{y} &= t^{-1}yt'.
\end{align*}
It is easy to check that $(\tilde{x}, \tilde{x}', \tilde{y}) \in \mathcal{S}_w$, hence $f_{t,t'}: \mathcal{S}_w \to \mathcal{S}_w$. It is also easy to check that $f_{t_1,t'_1} \circ f_{t_2,t'_2} = f_{t_1t_2, t'_1t'_2}$ for $(t_i, t'_i) \in H_w$, $i = 1, 2$, hence $f_{t,t'}$ defines an action of $H_w$ on $\mathcal{S}_w$. It is clear that this action extends the action \eqref{eq:6.6.1} of $T^F \times T'^F$ on $\mathcal{S}_w$. The theorem now follows from 6.5 and the

\begin{lemma}[6.7]
If the character $\theta \otimes \theta'$ of $T^F \times T'^F$ is trivial on $H_w \cap (T^F \times T'^F)$, then $\theta'^{-1} = \theta \circ \ad(F(w))$ (as characters of $Y(T')$).
\end{lemma}

The subgroup $H_w$ of $T \times T'$ is the kernel of the composite map
\[
T \times T' \xrightarrow{x^{-1}F(x), y^{-1}F(y)} T \times T' \xrightarrow{t, t' \mapsto \ad F(\dot{w})(t') \cdot t^{-1}} T.
\]
Applying the functor $Y$, we get that $Y(H_w)$ is the kernel $K$ of
\[
Y(T) \times Y(T') \xrightarrow{F-1} Y(T) \times Y(T') \xrightarrow{\ad F(w)(x') - x} Y(T).
\]
By (5.29), the triviality of $\theta \otimes \theta'$ on $H_w^\circ \cap (T^F \times T'^F)$ means that, when we view $\theta \otimes \theta'$ as a character of $Y(T) \times Y(T')$, it is trivial on $(F - 1)(K \otimes \mathbb{Q}) \cap (Y(T) \times Y(T'))$. Since the map $F - 1$ is injective, this intersection is
\[
\Ker(\ad F(w)(x') - x: Y(T) \times Y(T') \longrightarrow Y(T))
\]
and the assumption becomes the triviality of $\theta \otimes \theta'$ on $\Ker(\ad F(w)(x') - x)$, i.e., the identity $\theta \circ \ad F(w) \cdot \theta' = 1$.

We will now conjugate (6.2) with the character formula (4.2) to get quantitative results.

\begin{theorem}[6.8]
Let $\theta \in (T^F)^\wedge$, $\theta' \in (T'^F)^\wedge$. Then
\[
\langle R_T^\theta, R_{T'}^{\theta'} \rangle_{G^F} = \#\{w \in W(T, T')^F \mid \ad w(\theta') = \theta\}.
\]
\end{theorem}

\begin{theorem}[6.9]
\begin{equation}\label{eq:6.9.1}
\frac{1}{|G^F|} \sum_{\substack{u \in G^F \\ \text{unipotent}}} Q_{T,G}(u) Q_{T',G}(u) = \frac{|N(T, T')^F|}{|T^F| \cdot |T'^F|}.
\end{equation}
\end{theorem}

We will prove (6.8) and (6.9) simultaneously, by induction on the dimension of $G$.

\begin{lemma}[6.10]
If (6.9) holds, for $G$ replaced by $Z^0(s)$, where $s \in G^F$ is any semi-simple element not contained in the centre $Z$ of $G$, then
\begin{equation}\label{eq:6.10.1}
\langle R_T^\theta, R_{T'}^{\theta'} \rangle_{G^F} = \#\{w \in W(T, T')^F \mid \ad w(\theta') = \theta\} + \alpha
\end{equation}
where
\[
\alpha = \frac{1}{|G^F|} \sum_{\substack{s \in G^F \\ s \in T^F \cap Z \\ \text{semis.}}} \frac{1}{|Z^0(s)^F|} \sum_{\substack{g, g' \in G^F \\ g^{-1}sg \in T^F \\ g'^{-1}sg' \in T'^F}} \theta(g^{-1}sg)\overline{\theta'(g'^{-1}sg')} \sum_{\substack{u \in Z^0(s)^F \\ \text{unipotent}}} Q_{gTg^{-1}, Z^0(s)}(u) Q_{g'T'g'^{-1}, Z^0(s)}(u).
\]
In particular, (6.9) implies (6.8).
\end{lemma}

According to (4.2) we have:
\begin{align*}
\langle R_T^\theta, R_{T'}^{\theta'} \rangle &= \frac{1}{|G^F|} \sum_{g_1 \in G^F} \Tr(g_1, R_T^\theta) \overline{\Tr(g_1, R_{T'}^{\theta'})} \\
&= \frac{1}{|G^F|} \sum_{\substack{s \in G^F \\ \text{semis.}}} \frac{1}{|Z^0(s)^F|^2} \sum_{\substack{g, g' \in G^F \\ g^{-1}sg \in T^F \\ g'^{-1}sg' \in T'^F}} \theta(g^{-1}sg)\overline{\theta'(g'^{-1}sg')} \\
&\qquad \times \sum_{\substack{u \in Z^0(s)^F \\ \text{unipotent}}} Q_{gTg^{-1}, Z^0(s)}(u) Q_{g'T'g'^{-1}, Z^0(s)}(u).
\end{align*}
With our assumption, this can be written as
\begin{align*}
\frac{1}{|G^F|} \sum_{\substack{s \in G^F \\ s \notin Z \\ \text{semis.}}} \frac{1}{|Z^0(s)^F|} &\sum_{\substack{g, g' \in G^F \\ g^{-1}sg \in T^F \\ g'^{-1}sg' \in T'^F}} \theta(g^{-1}sg)\overline{\theta'(g'^{-1}sg')} \\
&\times \frac{|N_{Z^0(s)}(gTg^{-1}, g'T'g'^{-1})^F|}{|T^F| \cdot |T'^F|} + \alpha \\
&= \frac{1}{|G^F|} \sum_{\substack{s \in G^F \\ \text{semis.}}} \frac{1}{|Z^0(s)^F|} \sum_{\substack{g \in G^F \\ g^{-1}sg \in T^F}} \theta(g^{-1}sg) \\
&\qquad \times \sum_{\substack{n \in N_G(T, T')^F \\ n^{-1}g \in Z^0(s)^F}} \overline{\theta'(n^{-1}g^{-1}sgn)} \frac{|N_{Z^0(s)}(T, nT'n^{-1})^F|}{|T^F| \cdot |T'^F|} + \alpha.
\end{align*}
The formula $(g, g', n_1) \mapsto (g, n, n_1)$, $n = g'^{-1}n_1g$ establishes a one-to-one correspondence between the sets
\begin{align*}
&\{(g, g', n_1) \in G^F \times G^F \times G^F \mid g^{-1}sg \in T^F, g'^{-1}sg' \in T'^F, n_1 \in N_{Z^0(s)}(gTg^{-1}, g'T'g'^{-1})^F\} \\
&\{(g, n, n_1) \in G^F \times N_G(T, T')^F \times Z^0(s)^F \mid g^{-1}sg \in T^F\}
\end{align*}
It follows that
\begin{align*}
\langle R_T^\theta, R_{T'}^{\theta'} \rangle &= \frac{1}{|G^F|} \sum_{\substack{s \in G^F \\ \text{semis.}}} \frac{1}{|Z^0(s)^F|} \sum_{\substack{g \in G^F \\ n \in N_G(T, T')^F \\ g^{-1}sg \in T^F}} \theta(g^{-1}sg)\overline{\theta'(ng^{-1}sgn^{-1})} \\
&\qquad \times \frac{|Z^0(s)^F|}{|T^F| \cdot |T'^F|} + \alpha \\
&= \frac{1}{|G^F|} \cdot \frac{1}{|T^F| \cdot |T'^F|} \sum_{t \in T^F} \theta(t) \sum_{n \in N_G(T, T')^F} \overline{\theta'(ntn^{-1})} + \alpha \\
&= \alpha + \#\{w \in W(T, T')^F \mid \ad w(\theta') = \theta\}
\end{align*}
and (6.10.1) is proved.

\textit{Proof of 6.9.} Let $\bar{G} = G/Z$; the centre of $\bar{G}$ consists of the identity element. Both sides of \eqref{eq:6.9.1} remain unchanged when $G$ is replaced by $\bar{G}$; for the left hand side this follows from \eqref{eq:4.1.1}, while for the right hand side this is easy to check. Hence in order to prove (6.9), we may assume that the centre of $G$ has a single element; moreover we may assume by induction that the conclusion of (6.9) is true for $G$ replaced by $Z^0(s)$, where $s \in G^F$ is an arbitrary non-central semi-simple element, and for any two $F$-stable maximal tori in $Z^0(s)$. (To start the induction we may take $G$ to be the group with only one element in which case (6.9) is clear.) It follows from (6.3) that for any non-trivial character $\theta$ of $T^F$ one has $\langle R_T^1, R_{T'}^\theta \rangle = 0$. Substituting this information into (6.10), we see that if $T^F$ has any non-trivial character, then $\alpha = 0$ which proves (6.9). A similar argument applies when $|T^F| \neq 1$. It remains to check the special case where $|T^F| = |T'^F| = 1$. In this case, we must have $q = 2$, and $T$, $T'$ must be $\mathbb{F}_q$-split tori. In particular, $T$ and $T'$ are conjugate under $G^F$, hence
\[
\#\{w \in W(T, T')^F \mid \ad w(\theta') = \theta\} = |W(T)^F| = |W(T)|.
\]
Moreover $R_T^1$ and $R_{T'}^1$ are both equal to the representation of $G^F$ induced by the unit representation of $B^F$, where $B$ is an $F$-stable Borel subgroup of $G$. It follows that
\begin{equation}\label{eq:1.10}
\langle R_T^1, R_{T'}^1 \rangle = |B^F \backslash G^F / B^F| = |W(T)|
\end{equation}
Substituting in (6.10) we get (6.9) in this case. This completes the proof of (6.9) (and of (6.8) by (6.10)).

\section{Computations on semisimple elements}

The following result describes the Euler characteristic of the scheme $X_T \subset B$. Let $\sigma(G)$ (resp.\ $\sigma(T)$) be the $\mathbb{F}_q$-rank of $G$ (resp.\ $T$).

\begin{theorem}[7.1]
\[
\chi(X_T \subset B) = Q_{T,G}(e) = (-1)^{\sigma(G) - \sigma(T)} \frac{|G^F|}{\St_G(e) \cdot |T^F|}.
\]
(Here $\St_G$ is the Steinberg representation of $G^F$, identified with its character.)
\end{theorem}

We may assume that the centre of $G$ is reduced to a single element and that the statement is true when $G$ is replaced by $Z^0(s)$, where $s \in G^F$ is any non-central semi-simple element.

We first assume that $T^F$ has a non-trivial character $\theta$. The Steinberg representation occurs in $R_T^1 = \Ind_{B^F}^{G^F}(1)$ ($T^* \subset B^*$ as in 1.8), hence, by (6.3), it does not occur in $R_T^\theta$: $\langle R_T^\theta, \St_G \rangle = 0$. It is known that
\[
\St_G(g) = \begin{cases}
(-1)^{\sigma(G) - \sigma(Z^0(g))} \St_{Z^0(g)}(e), & g \in G^F \text{ semi-simple}, \\
0, & g \in G^F \text{ non-semi-simple}.
\end{cases}
\]
If we use 4.2, it follows:
\[
\frac{1}{|G^F|} \sum_{\substack{s \in G^F \\ \text{semis.}}} (-1)^{\sigma(G) - \sigma(Z^0(s))} \St_{Z^0(s)}(e) \sum_{\substack{g \in G^F \\ g^{-1}sg \in T^F}} Q_{gTg^{-1}, Z^0(s)}(e) \theta(g^{-1}sg) = 0
\]
(sum over the semi-simple elements $s$ of $G^F$).

By our assumption, we may substitute
\[
Q_{gTg^{-1}, Z^0(s)}(e) = (-1)^{\sigma(Z^0(s)) - \sigma(T)} \frac{|Z^0(s)^F|}{\St_{Z^0(s)}(e) \cdot |T^F|}
\]
for all $s \neq e$:
\[
(-1)^{\sigma(G)} \frac{1}{|G^F|} \sum_{\substack{s \in G^F \\ s \neq e \\ \text{semis.}}} \sum_{\substack{g \in G^F \\ g^{-1}sg \in T^F}} \theta(g^{-1}sg) + \St_G(e) Q_{T,G}(e) = 0.
\]
Since
\[
\frac{1}{|G^F|} \sum_{\substack{s \in G^F \\ \text{semis.}}} \sum_{\substack{g \in G^F \\ g^{-1}sg \in T^F}} \theta(g^{-1}sg) = \frac{1}{|T^F|} \sum_{t \in T^F} \theta(t) = 0
\]
since the character $\theta$ is non-trivial, $\sum_{t \in T^F} \theta(t) = 0$ and the desired formula for $Q_{T,G}(e)$ follows.

In the case where $|T^F| = 1$, $T$ must be an $\mathbb{F}_q$-split torus and $q = 2$. In this case, $Q_{T,G}(e) = |G^F|/|B^{*F}|$, where $B^*$ is any $F$-stable Borel subgroup (1.8). This agrees with the statement of the theorem and ends the proof.

\begin{corollary}[7.2]
For any semi-simple element $s \in G^F$,
\[
\Tr(s, R_T^\theta) = (-1)^{\sigma(Z^0(s)) - \sigma(T)} \frac{1}{\St_{Z^0(s)}(e)} \sum_{\substack{g \in G^F \\ g^{-1}sg \in T^F}} \theta(g^{-1}sg).
\]
\end{corollary}

This follows from (7.1) and 4.2. An equivalent statement is:

\begin{proposition}[7.3]
\[
(-1)^{\sigma(G) - \sigma(T)} R_T^\theta \otimes \St_G = \Ind_{T^F}^{G^F}(\theta).
\]
\end{proposition}

\begin{proposition}[7.4]
If $\theta$ is a non-singular character of $T^F$ (5.19) then $(-1)^{\sigma(G) - \sigma(T)} R_T^\theta$ can be represented by an actual $G^F$-module. If $\theta$ is in general position (5.15), $(-1)^{\sigma(G) - \sigma(T)} R_T^\theta$ is irreducible.
\end{proposition}

By embedding $G$ in a group $G_1$ with connected centre and the same derived group, as in (5.18), we reduce to the case where $\theta$ is in general position. It remains to use (6.8) and (7.1).

\begin{proposition}[7.5]
Let $s \in G^F$ be a semi-simple element. The character of
\begin{equation}\label{eq:7.5.1}
\sum_{\substack{T \ni s \\ (T)}} (-1)^{\sigma(G) - \sigma(T)} R_T^1
\end{equation}
equals $|Z(s)^F|$ on elements conjugate to $s$ in $G^F$ and zero on all other elements of $G^F$.
\end{proposition}

Let $\Phi_s$ be the character of \eqref{eq:7.5.1} and let $\Phi'_s$ be the class function on $G^F$ defined by
\[
\Phi'_s(g) = \begin{cases}
|Z(s)^F|, & g \text{ conjugate to } s \text{ in } G^F, \\
0, & \text{otherwise}.
\end{cases}
\]
In order to prove that $\Phi_s = \Phi'_s$ it is sufficient to prove that
\[
\langle \Phi_s, \Phi_s \rangle = \langle \Phi_s, \Phi'_s \rangle = \langle \Phi'_s, \Phi'_s \rangle \quad \text{(see (0.3))}.
\]
We have $\langle \Phi'_s, \Phi'_s \rangle = |Z(s)^F|$; from (6.8) we see that
\begin{align*}
\langle \Phi_s, \Phi_s \rangle &= \sum_{\substack{T \ni s \\ T' \ni s \\ (T), (T')}} \sum_{\substack{\theta \in (T^F)^\wedge \\ \theta' \in (T'^F)^\wedge \\ \theta(s^{-1})\theta'(s) \neq 0}} \#\{n \in N(T, T')^F \mid \ad n(\theta') = \theta\}/|T^F| \\
&= \frac{1}{|Z(s)^F|} \sum_{\substack{T \ni s \\ \theta \in (T^F)^\wedge}} \theta(s^{-1})\theta(nsn^{-1}) \quad (nsn^{-1} \in G^F) \\
&= \frac{1}{|Z(s)^F|} \sum_{\substack{T \ni s \\ n \in G^F \\ ns = sn}} 1 \\
&= \frac{1}{|Z(s)^F|} \cdot |Z(s)^F| \cdot \#\{F\text{-stable maximal tori in }Z^0(s)\} = |Z(s)^F|
\end{align*}
(by \cite[Cor.\ 14.16]{Stein}).

By (7.2), we have
\begin{align*}
\langle \Phi_s, \Phi'_s \rangle &= \frac{1}{|G^F|} \sum_{g \sim s} |Z(s)^F| \sum_{\substack{T \ni g \\ (T)}} (-1)^{\sigma(G) - \sigma(T)} R_T^1(g) \\
&= \frac{|Z(s)^F|}{|G^F|} \sum_{\substack{T \ni s \\ (T)}} \frac{(-1)^{\sigma(G) - \sigma(T)}}{\St_{Z^0(s)}(e)} \sum_{\substack{g \in G^F \\ g^{-1}sg \in T^F}} 1 \\
&= \frac{|Z(s)^F|}{\St_{Z^0(s)}(e)} \sum_{\substack{T \ni s \\ (T)}} \frac{(-1)^{\sigma(G) - \sigma(T)}}{|T^F|} \#\{g \in G^F \mid g^{-1}sg \in T^F\} \\
&= \frac{|Z(s)^F|}{\St_{Z^0(s)}(e)} \cdot \St_{Z^0(s)}(e) \cdot \#\{F\text{-stable maximal tori in }Z^0(s)\} = |Z(s)^F|
\end{align*}
and the proposition is proved.

\begin{corollary}[7.6]
For any $\rho \in \mathcal{R}(G^F)$ and any semi-simple element $s \in G^F$,
\begin{equation}\label{eq:7.6.1}
\Tr(s, \rho) = \frac{1}{\St_G(s)} \sum_{\substack{T \ni s \\ (T)}} (-1)^{\sigma(G) - \sigma(T)} \langle \rho, R_T^1 \rangle.
\end{equation}
In particular, if $s \in G^F$ and is regular semi-simple, contained in a unique torus $T$,
\begin{align}
\Tr(s, \rho) &= \frac{1}{|T^F|} \sum_{\theta \in (T^F)^\wedge} \theta(s) \langle \rho, R_T^\theta \rangle, \label{eq:7.6.2} \\
\dim \rho &= \frac{1}{\St_G(e)} \sum_{(T)} |T^F|(-1)^{\sigma(G) - \sigma(T)} \langle \rho, R_T^1 \rangle. \label{eq:7.6.3}
\end{align}
\end{corollary}

\begin{corollary}[7.7]
For any irreducible representation $\rho$ of $G^F$ there exists an $F$-stable maximal torus $T$ and a character $\theta$ of $T^F$ such that $\langle \rho, R_T^\theta \rangle \neq 0$.
\end{corollary}

This follows from \eqref{eq:7.6.3}.

\begin{definition}[7.8]
An irreducible representation $\rho$ of $G^F$ is said to be \textit{unipotent} if $\langle \rho, R_T^1 \rangle \neq 0$ for some $F$-stable maximal torus $T$.
\end{definition}

By (6.3) and (7.7), $\rho$ is unipotent if and only if $\langle \rho, R_T^\theta \rangle = 0$ for any $F$-stable maximal torus $T$ and any $\theta \in (T^F)^\wedge$, $\theta \neq 1$. For example, any irreducible component of $\Ind_{B^{*F}}^{G^F}(1)$ ($B^{*F}$ an $F$-stable Borel subgroup) is a unipotent representation. For unipotent representations, (7.6) becomes:

\begin{proposition}[7.9]
Let $\rho$ be a unipotent representation of $G^F$ and let $s \in G^F$ be a semi-simple element. Then
\[
\Tr(s, \rho) = \frac{1}{\St_G(s)} \sum_{\substack{T \ni s \\ (T)}} (-1)^{\sigma(G) - \sigma(T)} \langle \rho, R_T^1 \rangle.
\]
In particular, if $s \in G^F$ is regular semi-simple contained in a unique maximal torus $T$, $\Tr(s, \rho) = \langle \rho, R_T^1 \rangle$ and
\[
\dim \rho = \frac{1}{\St_G(e)} \sum_{(T)} |T^F| (-1)^{\sigma(G) - \sigma(T)} \langle \rho, R_T^1 \rangle.
\]
\end{proposition}

It follows that if $B$ is an $F$-stable Borel subgroup of $G$ and $T$ is an $F$-stable maximal torus in $B$, we have $\Tr(s, \rho) = \langle \rho, \Ind_{B^F}^{G^F}(1) \rangle$; in the case where $\langle \rho, \Ind_{B^F}^{G^F}(1) \rangle \neq 0$, this is a result of Curtis, Kilmoyer and Seitz (see C.\ W.\ Curtis, \textit{On the Values of Certain Irreducible Characters of Finite Chevalley Groups}, Ist.\ Naz.\ di Alta Mat.\ Symp.\ Mat.\ XIII, 1974, 343--355).

\begin{proposition}[7.10]
Let $G_{\ad}$ be the adjoint group of $G$. Then, the restriction to $G^F$ of a unipotent representation of $G_{\ad}^F$ is irreducible; non-isomorphic unipotent representations have non-isomorphic restrictions; and every unipotent representation of $G$ is such a restriction.
\end{proposition}

It suffices to check that, if $\sigma$, $\tau$ are unipotent representations of $G_{\ad}^F$, then $\langle \sigma, \tau \rangle_{G_{\ad}^F} = \langle \sigma, \tau \rangle_{G^F}$. Let $G_1$ be the image of $G^F$ in $G_{\ad}^F$. One has
\begin{align*}
\langle \sigma, \tau \rangle_{G^F} &= \langle \sigma, \tau \rangle_{G_1} = \langle \Ind_{G_1}^{G_{\ad}^F}(\Res_{G_1}^{G_{\ad}^F} \sigma), \tau \rangle_{G_{\ad}^F} \\
&= \sum_{\theta \in (G_{\ad}^F/G^F)^\wedge} \langle \sigma \otimes \theta, \tau \rangle_{G_{\ad}^F}.
\end{align*}
By the exclusion theorem 6.3 and (1.23), (1.27), $\langle \sigma \otimes \theta, \tau \rangle_{G_{\ad}^F} = 0$ for $\theta \neq 1$, hence the proposition.

\begin{proposition}[7.11]
Let $\rho$ be a virtual representation of $G^F$ such that $\Tr(su, \rho) = \Tr(s, \rho)$ for any $s \in G^F$ semi-simple, $u \in G^F$ unipotent, $su = us$. Let $T$ be an $F$-stable maximal torus and let $\theta$ be a character of $T^F$. Then
\[
\langle \rho, R_T^\theta \rangle_{G^F} = (-1)^{\sigma(G) - \sigma(T)} \langle \rho \otimes \St_G, R_T^\theta \rangle_{G^F} = \langle \rho, \theta \rangle_{T^F}.
\]
\end{proposition}

The Steinberg character is integral valued, hence by (7.3),
\begin{align*}
(-1)^{\sigma(G) - \sigma(T)} \langle \rho \otimes \St_G, R_T^\theta \rangle_{G^F} &= \langle \rho, (-1)^{\sigma(G) - \sigma(T)} \St_G \otimes R_T^\theta \rangle_{G^F} \\
&= \langle \rho, \Ind_{T^F}^{G^F}(\theta) \rangle_{G^F} = \langle \rho, \theta \rangle_{T^F}.
\end{align*}

We now prove the identity
\begin{equation}\label{eq:7.11.1}
\langle \rho, R_T^\theta \rangle_{G^F} = \langle \rho, \theta \rangle_{T^F}.
\end{equation}
The special case of this identity
\begin{equation}\label{eq:7.11.2}
\langle 1, R_T^\theta \rangle_{G^F} = 1
\end{equation}
follows from the fact that the Euler characteristic of $X_T \subset B/G^F$ (which is the same as the Euler characteristic of $\{g \in G \mid g^{-1}Fg \in FU\}/G^F \times T^F = FU/T^F$ by (1.17)) equals $1$. If we use (4.2), \eqref{eq:7.11.2} can be written in the form
\begin{equation}\label{eq:7.11.3}
\frac{1}{|G^F|} \sum_{\substack{s \in G^F \\ \text{semis.}}} \frac{1}{|Z^0(s)^F|} \sum_{\substack{g \in G^F \\ g^{-1}sg \in T^F}} \sum_{\substack{u \in Z^0(s)^F \\ \text{unipotent}}} Q_{gTg^{-1}, Z^0(s)}(u) = 1.
\end{equation}
This implies, by induction on the dimension of $G$, that
\begin{equation}\label{eq:7.11.4}
\frac{1}{|G^F|} \sum_{\substack{u \in G^F \\ \text{unipotent}}} Q_{T,G}(u) = \frac{1}{|T^F|}.
\end{equation}
Indeed, by substituting \eqref{eq:7.11.4} in \eqref{eq:7.11.3} we find a true identity:
\[
\frac{1}{|G^F|} \sum_{\substack{s \in G^F \\ \text{semis.}}} \frac{|G^F|}{|Z^0(s)^F|} \cdot \frac{1}{|T^F|} \cdot \#\{g \in G^F \mid g^{-1}sg \in T^F\} = 1.
\]
We apply (4.2) again and use \eqref{eq:7.11.4} to compute:
\begin{align*}
\langle \rho, R_T^\theta \rangle_{G^F} &= \frac{1}{|G^F|} \sum_{\substack{s \in G^F \\ \text{semis.}}} \frac{1}{|Z^0(s)^F|} \sum_{\substack{g \in G^F \\ g^{-1}sg \in T^F}} \sum_{\substack{u \in Z^0(s)^F \\ \text{unipotent}}} Q_{gTg^{-1}, Z^0(s)}(u) \theta(g^{-1}sg) \Tr(s, \rho) \\
&= \frac{1}{|G^F|} \sum_{\substack{s \in G^F \\ \text{semis.}}} \sum_{\substack{g \in G^F \\ g^{-1}sg \in T^F}} \frac{1}{|T^F|} \theta(g^{-1}sg) \Tr(s, \rho) \\
&= \frac{1}{|T^F|} \sum_{t \in T^F} \theta(t) \Tr(t, \rho) = \langle \rho, \theta \rangle_{T^F}.
\end{align*}

\begin{proposition}[7.12]
Let $\rho$ be as in (7.11). Then
\begin{align}
\label{eq:7.12.1} \rho &= \sum_{(T)} \frac{(-1)^{\sigma(G) - \sigma(T)}}{|W(T)^F|} \sum_{\theta \in (T^F)^\wedge} \langle \rho, \theta \rangle_{T^F} R_T^\theta, \\
\label{eq:7.12.2} \rho \otimes \St_G &= \sum_{(T)} \frac{(-1)^{\sigma(G) - \sigma(T)}}{|W(T)^F|} \sum_{\theta \in (T^F)^\wedge} \langle \rho, \theta \rangle_{T^F} R_T^\theta,
\end{align}
where $(T)$ means sum over all $G^F$-conjugacy classes of $F$-stable maximal tori $T$.
\end{proposition}

The character of $\rho \otimes \St_G$ is a linear combination of the form $\sum_{(T, \theta)} c_{T, \theta} R_T^\theta$ (sum over all $G^F$-conjugacy classes of pairs $(T, \theta)$). The coefficients $c_{T, \theta}$ are determined by (7.11) and \eqref{eq:7.12.2} follows.

Now let $\rho'$ (resp.\ $\rho''$) be the right hand side of \eqref{eq:7.12.1} (resp.\ \eqref{eq:7.12.2}). We have clearly
\begin{equation}\label{eq:7.12.3}
\langle \rho', \rho' \rangle = \langle \rho'', \rho'' \rangle
\end{equation}
and by (7.11),
\begin{equation}\label{eq:7.12.4}
\langle \rho, \rho' \rangle = \langle \rho \otimes \St_G, \rho'' \rangle
\end{equation}
We note also that
\begin{equation}\label{eq:7.12.5}
\langle \rho, \rho \rangle = \langle \rho \otimes \St_G, \rho \otimes \St_G \rangle
\end{equation}
This follows from the fact that the number of unipotents in $Z^0(s)^F$ equals $\St_G(s)^2$ for any semi-simple element $s \in G^F$ (see \cite[Thm.\ 15.1]{Stein}). By \eqref{eq:7.12.2} we have
\[
\langle \rho \otimes \St_G - \rho', \rho \otimes \St_G - \rho' \rangle = 0.
\]
This, together with \eqref{eq:7.12.3}, \eqref{eq:7.12.4} and \eqref{eq:7.12.5} implies:
\[
\langle \rho - \rho', \rho - \rho' \rangle = 0
\]
hence
\[
\rho = \rho'.
\]

\begin{remark}[7.13]
It is known that
\[
\langle \rho, \rho \rangle_{G^F} = \sum_{(T)} \frac{1}{|T^F|} \langle \rho, \rho \rangle_{T^F}
\]
(N.\ Kawanaka, \textit{A theorem on finite Chevalley groups}, Osaka J.\ Math.\ 10 (1973), 1--13); this could be also deduced from \eqref{eq:7.12.1} and (6.8).
\end{remark}

\begin{corollary}[7.14]
\begin{align}
\label{eq:7.14.1} 1 &= \sum_{(T)} \frac{(-1)^{\sigma(G) - \sigma(T)}}{|W(T)^F|} R_T^1, \\
\label{eq:7.14.2} \St_G &= \sum_{(T)} \frac{(-1)^{\sigma(G) - \sigma(T)}}{|W(T)^F|} R_T^1.
\end{align}
In particular,
\[
\langle \St_G, R_T^1 \rangle = (-1)^{\sigma(G) - \sigma(T)}.
\]
\end{corollary}

Here is an alternative proof of \eqref{eq:7.14.1}. In the language of (1.4), \eqref{eq:7.14.1} asserts that
\[
\sum_{w \in W} R^1(w) = |W| \cdot 1.
\]
Since $(X(w))_{w \in W}$ is a partition of the flag manifold into locally closed subschemes, we have
\[
\sum_{w \in W} R^1(w) = \sum_i (-1)^i H_c^i(X_G)
\]
as virtual representations of $G^F$. The action of $G^F$ on $X_G$ is the restriction of the action of the connected group $G$; by (6.5), $G^F$ acts trivially on $H_c^*(X_G)$ and it remains to use the fact that the Euler characteristic of $X_G$ equals $|W|$.

\begin{corollary}[7.15]
\begin{align}
\label{eq:7.15.1} \St_G &= \sum_{(T)} \frac{(-1)^{\sigma(G) - \sigma(T)}}{|W(T)^F|} \Ind_{T^F}^{G^F}(1), \\
\label{eq:7.15.2} \St_G \otimes \St_G &= \sum_{(T)} \frac{1}{|W(T)^F|} \Ind_{T^F}^{G^F}(1).
\end{align}
\end{corollary}

This follows from (7.1) by tensoring with $\St_G$, and using (7.3). The formula \eqref{eq:7.15.1} is due to B.\ Srinivasan (\textit{On the Steinberg character of a finite simple group of Lie type}, J.\ Australian Math.\ Soc.\ 12 (1971), 1--14).

\section{Induced and cuspidal representations}

\subsection{8.1} Let $P$ be an $F$-stable parabolic subgroup of $G$ and let $T \subset P$ be an $F$-stable maximal torus. We denote by $U_P$ the unipotent radical of $P$. The quotient group $P/U_P$ is a connected reductive algebraic group acted on by $F$. Let $\pi: P \to P/U_P$ be the canonical projection. $\pi$ induces an isomorphism $T \xrightarrow{\sim} \pi(T)$ hence also an isomorphism $T^F \xrightarrow{\sim} \pi(T)^F$. Let $\theta$ be a character of $T^F$ and let $\tilde{\theta}$ be the corresponding character of $\pi(T)^F$. We denote by $R_T^\theta \cdot P$ the image of the virtual representation $R_{\pi(T)}^{\tilde{\theta}}$ of $(P/U_P)^F$ under the canonical embedding $\mathcal{R}((P/U_P)^F) \subset \mathcal{R}(P^F)$. With these notations, we have the following generalisation of 1.10:

\begin{proposition}[8.2]
\[
R_T^\theta \cdot P = \Ind_{P^F}^{G^F}(R_T^\theta \cdot P).
\]
\end{proposition}

We choose a Borel subgroup $B$ in $G$ such that $T \subset B \subset P$. Let $\mathcal{P}$ be the set of all parabolic subgroups $P'$ in $G$ such that $P'$ and $P$ are conjugate under $G^F$; this is a finite set. We have
\[
\dot{X}_T \subset B = \coprod_{P' \in \mathcal{P}} \dot{X}(P')
\]
where
\[
\dot{X}(P') = \{g \in G \mid g^{-1}Fg \in FU, \ gPg^{-1} = P'\}/U \cap FU,
\]
with $U$ the unipotent radical of $B$. (Note that $g^{-1}F(g) \in FU$ implies $F(gPg^{-1}) = gPg^{-1}$.) Let $P' \in \mathcal{P}$ and let $g_1 \in G^F$ be such that $g_1Pg_1^{-1} = P'$. If $g \in G$, $g^{-1}F(g) \in FU$, $gPg^{-1} = P'$, we have
\[
g_1^{-1}g \in P \quad \text{and} \quad (g_1^{-1}g)^{-1}F(g_1^{-1}g) \in FU;
\]
by sending $g$ to $\pi(g_1^{-1}g)$ we get an isomorphism $\dot{X}(P') \xrightarrow{\sim} \dot{X}_{\pi(T) \subset \pi(B)}$. It follows easily that the $G^F$-module $H_c^*(\dot{X}_T \subset B)$ is canonically isomorphic to the $G^F$-module induced by $H_c^*(\dot{X}_{\pi(T) \subset \pi(B)})$ regarded as a $P^F$-module; moreover this isomorphism is compatible with the action of $T^F$. The proposition follows.

\begin{theorem}[8.3]
Let $T$ be an $F$-stable maximal torus in $G$ such that $T$ is not contained in any $F$-stable, proper parabolic subgroup of $G$. Let $\theta$ be a non-singular character of $T^F$ (cf.\ 5.19). Then $(-1)^{\sigma(G) - \sigma(T)} R_T^\theta$ can be represented by a cuspidal $G^F$-module.
\end{theorem}

Applying (7.5) for $s$ equal to the identity element of $G^F$, we see that the regular representation of $G^F$ is equal to:
\[
\Ind_{B^F}^{G^F}(1) = \frac{1}{\St_G(e)} \sum_{(T)} \sum_{\theta \in (T^F)^\wedge} (-1)^{\sigma(G) - \sigma(T)} R_T^\theta.
\]
We apply this formula with $G$ replaced by $P/U_P$ ($P$ as in (8.1)); we then regard the regular representation of $(P/U_P)^F$ as a representation of $P^F$ on which $U_P$ acts trivially (i.e., as $\Ind_{U_P^F}^{P^F}(1)$) and we induce it to $G^F$:
\[
\Ind_{U_P^F}^{G^F}(1) = \Ind_{P^F}^{G^F}(\Ind_{U_P^F}^{P^F}(1)) = \Ind_{P^F}^{G^F}\left(\frac{1}{\St_{P/U_P}(e)} \sum_{(T')} (-1)^{\sigma(P/U_P) - \sigma(T')} R_{T'}^{1'}\right)
\]
where $R_{T'}^{1'}$ is regarded as an element in $\mathcal{R}(P^F)$ in the natural way. We now use (8.2) and the fact that any $F$-stable maximal torus $T'$ in $P/U_P$ can be lifted to an $F$-stable maximal torus $T$ in $P$ in precisely $|U_P^F|$ different ways:
\begin{equation}\label{eq:8.3.1}
\Ind_{U_P^F}^{G^F}(1) = \frac{1}{\St_G(e)} \sum_{\substack{(T) \\ T \subset P}} |U_P^F| \cdot (-1)^{\sigma(G) - \sigma(T)} R_T^1.
\end{equation}
If $P \neq G$ and $T$ is as in the statement of the theorem, we see from \eqref{eq:8.3.1} and (6.8) that
\begin{equation}\label{eq:8.3.2}
\langle (-1)^{\sigma(G) - \sigma(T)} R_T^\theta, \Ind_{P^F}^{G^F}(1) \rangle = 0 \quad \text{for any } \theta \in (T^F)^\wedge.
\end{equation}
If $\theta$ is non-singular, $(-1)^{\sigma(G) - \sigma(T)} R_T^\theta$ can be represented by an actual $G^F$-module (7.4). This $G^F$-module is cuspidal by \eqref{eq:8.3.2}.

\section{A vanishing theorem}

\subsection{9.1} Let $G$ be a reductive algebraic group over $k$. Let $w = s_1 \cdots s_n$ be a minimal expression for an element of the Weyl group $W$ of $G$ (0.4). The following desingularization of the closure $\overline{O(w)}$ of $O(w) \subset X \times X$ (1.2) has been considered by H.\ C.\ Hansen \cite{Han} and M.\ Demazure \cite{Dem}.

\begin{definition}[9.2]
$\overline{O}(s_1, \dots, s_n)$ is the space of sequences $(B_0, \dots, B_n)$ of Borel subgroups, with $B_{i-1}$ and $B_i$ in relative position $e$ or $s_i$.
\end{definition}

The maps
\[
\overline{O}(s_1, \dots, s_i) \longrightarrow \overline{O}(s_1, \dots, s_{i-1})
\]
express $\overline{O} = \overline{O}(s_1, \dots, s_n)$ as an iterated fibre space over $X$, with fibre $\mathbb{P}^1$ (1.2). The subspace $D_i \subset \overline{O}$ where $B_{i-1} = B_i$ is the inverse image of a section of $\overline{O}(s_1, \dots, s_i) \to \overline{O}(s_1, \dots, s_{i-1})$; it is a smooth divisor, and $D = \bigcup D_i$ is a divisor with normal crossings. The map $(B_0, \dots, B_n) \mapsto (B_0, B_n)$ induces an isomorphism $\overline{O} - D \xrightarrow{\sim} O(w)$, and hence maps $\overline{O}$ to $\overline{O(w)}$; $\overline{O}$ is a resolution of singularities for $\overline{O(w)}$.

\subsection{9.3} Let $T^*$, $B^*$ and $U^*$ be as in (1.7). For any character $\chi: T^* \to \mathbb{G}_m$, the $T^*$-torsor $\dot{E}$ over $X$ defined in (1.7) gives rise to a line bundle $\dot{E}_\chi$, provided with $\chi: \dot{E} \to \dot{E}_\chi$ such that $\chi(et) = \chi(e)\chi(t)$. For $\dot{w} \in N(T^*)$ with image $w$ in $W = N(T^*)/T^*$, the $w$-map of $T^*$-torsors over $O(w) \subset X \times X$ constructed in (1.7) induces an isomorphism
\[
T(\dot{w}): \pr_2^* \dot{E}_\chi \xrightarrow{\sim} \pr_1^* \dot{E} \otimes \ad\dot{w} = \pr_1^* \dot{E}_{w^{-1}(\chi)}
\]
which makes the following diagram commute:
\[
\begin{array}{ccc}
\pr_2^* \dot{E} & \xrightarrow{\dot{w}} & \pr_1^* \dot{E} \\
\pr_2^* \chi \downarrow & & \downarrow \pr_1^* \chi \\
\pr_2^* \dot{E}_\chi & \xrightarrow{T(\dot{w})} & \pr_1^* \dot{E}_{w^{-1}(\chi)}.
\end{array}
\]
We will investigate its behaviour at infinity.

\subsection{9.4} Let $X(T^*)$ be the character group of $T^*$ and let $C \subset X(T^*) \otimes \mathbb{R}$ be the fundamental chamber (corresponding to $B^*$). If $C_1$ and $C_2$ are two chambers, we write $D(C_1, C_2)$ for the intersection of the (closed) radicial half spaces containing both $C_1$ and $C_2$, and $D^0(C_1, C_2)$ for its interior.

\begin{proposition}[9.5]
Let $\overline{O(w)}'$ be the normalization of the closure $\overline{O(w)}$ of $O(w)$ in $X \times X$. The map $T(\dot{w}): \pr_2^* \dot{E}_\chi \to \pr_1^* \dot{E}_{w^{-1}(\chi)}$ extends over $\overline{O(w)}'$ if and only if $\chi \in D(C, -wC)$. It vanishes outside of $O(w)$ if and only if $\chi \in D^0(C, -wC)$.
\end{proposition}

Let us consider $T(\dot{w})$ as a rational section of $\pr_2^* \dot{E}_{\chi^{-1}} \otimes \pr_1^* \dot{E}_{w^{-1}(\chi)}$; it suffices to prove that, with the notations of (9.2), for each $i$, the order $\nu_i$ of $T(\dot{w})$ along the divisor $D_i \subset \overline{O}(s_1, \dots, s_n)$ is $\ge 0$ for $\chi \in D(C, -wC)$, $> 0$ for $\chi \in D^0(C, -wC)$. The convex region $D(C, -wC)$ is the intersection of the radicial half spaces containing $C$ but not $wC$. Put $w_i = s_1 \cdots s_i$. The gallery $(C, w_1C, w_2C, \dots, w_nC)$ connects $C$ to $wC$; its walls are all the radicial hyperplanes that separate $C$ and $wC$. The wall between $w_{i-1}C$ and $w_iC$ is defined by the root $w_{i-1}(\alpha_i)$, where $\alpha_i$ is the simple root corresponding to $s_i$; $C$ and $w_{i-1}C$ are on the same side of it. The convex region $D(C, -wC)$ is hence defined by the inequalities
\[
\langle \chi, H_{w_{i-1}(\alpha_i)} \rangle \ge 0.
\]
(Recall that $H_\alpha$ denotes the coroot corresponding to a root $\alpha$, cf.\ (5.1).) We will prove that
\begin{equation}\label{eq:9.5.1}
\nu_i = \langle \chi, H_{w_{i-1}(\alpha_i)} \rangle.
\end{equation}
Let us lift the decomposition $w = (s_1 \cdots s_{i-1}) s_i (s_{i+1} \cdots s_n)$ into a decomposition $\dot{w} = \dot{w}' \dot{s}_i \dot{w}''$ in $N(T^*)$. We have
\[
T(\dot{w}) = T(\dot{w}'') \circ T(\dot{s}_i) \circ T(\dot{w}').
\]
When we approach a general point of $D_i$, the isomorphisms $T(\dot{w}')$ and $T(\dot{w}'')$ extend, and we are left to prove that the order along $D_i$ of
\[
T(\dot{s}_i): \dot{E}_{w'^{-1}(\chi)} \longrightarrow \dot{E}_{s_i w'^{-1}(\chi)}
\]
is
\[
\langle \chi, H_{w_{i-1}(\alpha_i)} \rangle = \langle w'^{-1}(\chi), H_{\alpha_i} \rangle;
\]
we are reduced to the case where $w$ is a fundamental reflection: all computations occur within a minimal parabolic subgroup $P$, or $P/U_P$, or the derived group of $P/U_P$, or its universal covering $SL(2)$. Let us make a direct check for $G = SL(2)$, $w \neq e$ and $\chi$ the fundamental weight.

(a) We take $SL(2)$ in its obvious representation $k^2$, $T^* =$ diagonal matrices, $B^* =$ upper triangular matrices. The weight $\chi$ is
\[
\begin{pmatrix} a & 0 \\ 0 & a^{-1} \end{pmatrix} \mapsto a, \quad \text{and} \quad \langle \chi, H_\alpha \rangle = 1.
\]
(b) $X = \mathbb{P}^1$, the space of homogeneous lines of $k^2$ and the fibre of $\dot{E}_\chi$ at $x$ is the corresponding line.

(c) Up to a scalar, $T(\dot{w})$ associates to $u \in (\dot{E}_\chi)_x$ the linear form $u \wedge v$ on $(\dot{E}_\chi)_y$, for $x \neq y$. It vanishes simply for $y \to x$.

\subsection{9.6} The assumptions (0.1.1) are in force from now on. Under the assumptions of (1.8), we have
\[
F^* \dot{E}_\chi = \dot{E}_{F^*\chi} \quad (\text{where } F^*\chi = \chi \circ F).
\]
On the inverse image $X(w)'$ of the graph of the Frobenius map $F: X \to X$ in $\overline{O(w)}'$, the line bundle $\pr_2^* \dot{E}_\chi^{-1} \otimes \pr_1^* \dot{E}_{w^{-1}(\chi)}$ is hence isomorphic to
\[
\pr_1^*(\dot{E}_{Fw^{-1}(\chi) - \chi}^{-1}).
\]
It will be ample if $\dot{E}_{Fw^{-1}(\chi) - \chi}$ is ample on $X$, i.e., if $F^*w^{-1}\chi - \chi \in -C$. On the other hand, if $\chi \in D^0(C, -wC)$, the section $T(\dot{w})$ of it has as zero set the complement of $X(w)$ in the projective variety $X(w)'$. If both conditions can be simultaneously fulfilled, then $X(w)$ is affine. In terms of $\mu = -w^{-1}\chi$, the conditions read
\begin{align}
\label{eq:9.6.1} \mu &\in D^0(C, -w^{-1}C), \\
\label{eq:9.6.2} F^*\mu - w\mu &\in -C.
\end{align}

\begin{theorem}[9.7]
If there is $\mu \in X(T^*) \otimes \mathbb{R}$ satisfying \eqref{eq:9.6.1} and \eqref{eq:9.6.2} then $X(w)$ is affine. As a consequence, $X(w)$ is affine as soon as $q$ is larger than the Coxeter number $h$ of $G$.
\end{theorem}

It remains to prove that if $q > h$, then some $\mu$ satisfies \eqref{eq:9.6.1} and \eqref{eq:9.6.2}. The first condition will be fulfilled if $\mu \in C^0$, that is, if $\langle \mu, H_\alpha \rangle > 0$ for each simple root $\alpha$. We take $\mu$ so that $\langle \mu, H_\alpha \rangle = 1$ for each simple root $\alpha$. We then have $\langle F\mu, H_\alpha \rangle = q$, and $\langle w\mu, H_\alpha \rangle = \langle \mu, w^{-1}H_\alpha \rangle$. If $H = \sum n_\alpha H_\alpha$ is the highest coroot, $\sum n_\alpha = h - 1$ and
\[
\langle F\mu - w\mu, H_\alpha \rangle \ge q - \langle \mu, H \rangle = q - h + 1 > 0,
\]
hence \eqref{eq:9.6.2}.

\begin{theorem}[9.8]
If the character $\theta: T(w)^F \to \overline{\mathbb{Q}}_\ell^*$ is non-singular, then
\[
H_c^*(\dot{X}(\dot{w}), \overline{\mathbb{Q}}_\ell)_\theta = H^*(\dot{X}(\dot{w}), \overline{\mathbb{Q}}_\ell)_\theta.
\]
\end{theorem}

\begin{corollary}[9.9]
If $X(w)$ is affine and $\theta$ is non-singular, then
\[
H_c^i(\dot{X}(\dot{w}), \overline{\mathbb{Q}}_\ell)_\theta = 0 \quad \text{for } i \neq l(w).
\]
\end{corollary}

Let us deduce (9.9) from (9.8). With the notations of (1.9), (9.8) means that
\begin{equation}\label{eq:9.9.1}
H_c^*(X(w), \mathcal{F}_\theta) = H^*(X(w), \mathcal{F}_\theta).
\end{equation}
If $X(w)$ is affine, then $H^i(X(w), \mathcal{F}_\theta) = 0$ for $i > \dim X(w) = l(w)$ (\cite[XIV, 3.2]{SGA4}). The sheaf $\mathcal{F}_\theta$ is locally constant, its dual is $\mathcal{F}_{\theta^{-1}}$, and $X(w)$ is smooth and purely of dimension $l(w)$. By Poincar\'e duality, $H_c^i(X(w), \mathcal{F}_\theta)$ is dual (up to a twist) to $H^{2l(w)-i}(X(w), \mathcal{F}_{\theta^{-1}})$, hence vanishes for $2l(w) - i > l(w)$, i.e., for $i < l(w)$. This proves the corollary.

\subsection{9.10} We first construct a nice compactification of $X(w)$. Let $w = s_1 \cdots s_n$ be a minimal expression for $w$. We define $\overline{X}(s_1, \dots, s_n)$ to be the space of sequences $(B_0, \dots, B_n)$ of Borel subgroups of $G$, with $B_n = FB_0$ and with $B_{i-1}$ and $B_i$ in relative positions $s_i$ or $e$. It is the inverse image of the graph of Frobenius by the map
\[
\overline{O}(s_1, \dots, s_n) \xrightarrow{(B_0, B_n)} X \times X.
\]
In other words, if $\Gamma_F \hookrightarrow X \times X$ is the inclusion in $X \times X$ of the graph of Frobenius, it is the fibre product $\overline{O} \times_{X \times X} \Gamma_F$:
\[
\begin{array}{ccc}
\overline{X}(s_1, \dots, s_n) & \subseteq & \Gamma_F \\
\downarrow & & \downarrow \\
\overline{O}(s_1, \dots, s_n) & \xrightarrow{(B_0, B_n)} & X \times X.
\end{array}
\]

\begin{lemma}[9.11]
$\Gamma_F$ is transverse to $\overline{O}(s_1, \dots, s_n)$, as well as to any intersection of the divisors $D_i$. The fibre product $\overline{X}(s_1, \dots, s_n)$ is hence a smooth compactification of $X(w)$, with a divisor with normal crossings (sum of the traces $\bar{D}_i$ of the $D_i$) at infinity.
\end{lemma}

We will show that $F$ is transverse to any smooth $G$-equivariant scheme $\pi: Y \to X \times X$ over $X \times X$ (with $G$ acting diagonally on $X \times X$); that is, if $\pi(y) \in \Gamma_F$, the sum of the tangent space to $\Gamma_F$ at $\pi(y)$ and of the image of $d\pi$ is the whole tangent space of $X \times X$ at $\pi(y)$. Indeed

(a) the tangent space of $\Gamma_F$ at $\pi(y)$ is $T_x \times \{0\}$;

(b) the image of $d\pi$ contains the image of $\Lie(G)$ by the derivative at $e \in G$ of $g \mapsto g \cdot y$ (by the equivariance of $Y$). Since the space $X$ is homogeneous with reduced stabilizers this image projects onto $T_x$ by the second projection.

\subsection{9.12} We now investigate how the covering $\dot{X}(\dot{w})$ (with structural group $T(w)^F$) ramifies along the divisor $\bar{D}_i$. The structural group being of order prime to $p$, the ramification is tame. On each connected component of $\bar{D}_i$, it gives rise to a homomorphism $\tau_i: \mathbb{Z}(1) \to T(w)^F$.

Let $x(t) \in \overline{X}(s_1, \dots, s_n)$ be a one parameter family with $x(0)$ a point of $\bar{D}_i$ (not on any $\bar{D}_j$, $j \neq i$), and with the tangent vector $\dot{x}(0)$ transverse to $\bar{D}_i$. In technical terms: for $S = \text{spectrum of the Henselization of }k[t]\text{ at }(t)$, $x$ is a morphism $x: S \to \overline{X}(s_1, \dots, s_n)$ such that the closed point $t = 0$ is mapped to a point of $\bar{D}_i$ not on any $\bar{D}_j$, $j \neq i$, and that the inverse image of $\bar{D}_i$ is the reduced scheme $t = 0$. Put $x(t) = (B_0(t), \dots, B_n(t))$ and let $E_i$ be the pull back by $B_i$ of the $T$-torsor $\dot{E}$ on $X$ (1.7). If we factor $\dot{w}$ as in 9.5, the isomorphism $T(\dot{w}): E_0 \to E_n$ (over the generic point of $S$) factors as $T(\dot{w}) = T(\dot{w}'') \circ T(\dot{s}_i) \circ T(\dot{w}')$ and $T(\dot{w}')$ and $T(\dot{w}'')$ extend over $S$; $T(\dot{s}_i)$ doesn't; rather, the composite $x^*T(\dot{s}_i)(x^*H_{w_{i-1}(\alpha_i)}(t))$ does (9.5), hence $T(\dot{w})$ is of the form
\[
T(\dot{w})(x) = T_0(x) \circ H_{w_{i-1}(\alpha_i)}(\alpha_i(t))
\]
where $T_0: E_0 \to E_n$ extends over $S$. The pull back under $x$ of the $T(w)^F$-torsor $\dot{X}(\dot{w})$, to the generic point of $S$, is given by the equation
\begin{equation}\label{eq:9.12.1}
F(u) = T_0(u \cdot H_{w_{i-1}(\alpha_i)}(\alpha_i(t))).
\end{equation}
Let $u_0$ be a solution, over $S$, of the equation
\[
F(u_0) = T_0(u_0)
\]
(this exists, because $S$ is strictly Henselian and $F$ and $T_0$ extend over $S$); if we put $u = u_0y$, \eqref{eq:9.12.1} becomes
\begin{equation}\label{eq:9.12.2}
y^{-1} \ad\dot{w}F(y) = H_{w_{i-1}(\alpha_i)}(\alpha_i(t)).
\end{equation}
In other words:

\begin{lemma}[9.13]
The $T(w)^F$-torsor $\dot{X}(\dot{w})$ ramifies along $\bar{D}_i$ in the same way as the pull back under $H_{w_{i-1}(\alpha_i)}: \mathbb{G}_m \to T$ of the Lang covering of $T(w)$ ramifies at $0$.
\end{lemma}

The tame fundamental group of $\mathbb{G}_m$ is $\hat{\mathbb{Z}}_{p'}(1)$, that of $T$ is $Y \otimes \hat{\mathbb{Z}}_{p'}(1)$ (for $Y$ the dual of the character group of $T$) and $H_{w_{i-1}(\alpha_i)}$ induces $x \mapsto x \otimes H_{w_{i-1}(\alpha_i)}: \hat{\mathbb{Z}}_{p'}(1) \to Y \otimes \hat{\mathbb{Z}}_{p'}(1)$. The Lang covering of $T(w)^F$ gives rise to the map $Y \otimes \hat{\mathbb{Z}}_{p'}(1) \to T(w)^F$ described in (5.2). Hence the pull back (9.13) corresponds to the composite
\[
\hat{\mathbb{Z}}_{p'}(1) \xrightarrow{x \mapsto x \otimes H_{w_{i-1}(\alpha_i)}} Y \otimes \hat{\mathbb{Z}}_{p'}(1) \longrightarrow T(w)^F.
\]
For $\theta$ a character of $T(w)^F$, $\mathcal{F}_\theta$ will ramify along $\bar{D}_i$ if and only if $\theta$ is not trivial on $H_{w_{i-1}(\alpha_i)} \hat{\mathbb{Z}}_{p'}(1)$. In that case, all higher direct images $R^i j_* \mathcal{F}_\theta$ of $j_! \mathcal{F}_\theta$ ($i > 0$) under $j: X(w) \hookrightarrow \overline{X}(s_1, \dots, s_n)$ will vanish on $\bar{D}_i$. In particular, we have

\begin{lemma}[9.14]
If $\theta$ is non-singular, $j_! \mathcal{F}_\theta = j_* \mathcal{F}_\theta$ (extension by zero) and $R^i j_* \mathcal{F}_\theta = 0$ for $i > 0$.
\end{lemma}

\textit{Proof of 9.8} (in the guise \eqref{eq:9.9.1}). By (9.14), the Leray spectral sequence for $j$ reads
\[
H^*(\overline{X}(s_1, \dots, s_n), j_! \mathcal{F}_\theta) \xrightarrow{\sim} H^*(X(w), \mathcal{F}_\theta).
\]
The space $\overline{X}(s_1, \dots, s_n)$ being a compactification of $X(w)$, the left hand side is by definition
\[
H_c^*(X(w), \mathcal{F}_\theta).
\]

\begin{remark}[9.15.1]
By using arguments parallel to those of (1.6), one can show that the conclusion of (9.9): ``for $\theta$ non-singular, $H_c^i(\dot{X}(\dot{w}), \overline{\mathbb{Q}}_\ell)_\theta = 0$ for $i \neq l(w)$'' holds true as soon as there is $w'$ in the $F$-conjugacy class of $w$ such that $X(w')$ is affine. The criterion (9.7) can be used to check that this is always the case for the classical groups and for $G_2$ except the case $(G_2, q = 2, w = \text{Coxeter})$ for which a different method applies.
\end{remark}

\begin{remark}[9.15.2]
Let $(T', \theta')$ be maximally split (5.25); under an isomorphism $T(w) \xrightarrow{\sim} T'$ (inducing $T(w)^F \xrightarrow{\sim} T'^F$), as in (1.18), $\theta'$ becomes a character of $T(w)^F$.
\end{remark}

\begin{thebibliography}{99}

\bibitem{Bou} N.\ Bourbaki, \textit{Groupes et Alg\`ebres de Lie}, Chap.\ IV, V, VI, Hermann, Paris, 1968.

\bibitem{ChRee} B.\ Chang, R.\ Ree, \textit{The characters of $G_2(q)$}, Symp.\ Math.\ XIII (1974), 395--413.

\bibitem{Dem} M.\ Demazure, \textit{D\'esingularisation des vari\'et\'es de Schubert g\'en\'eralis\'ees}, Ann.\ Sci.\ Ec.\ Norm.\ Sup.\ 7 (1974), 53--88.

\bibitem{DL1} P.\ Deligne, G.\ Lusztig, \textit{Representations of reductive groups over finite fields}, Proc.\ Int.\ Cong.\ Math.\ Vancouver 1974, 355--356.

\bibitem{DL2} P.\ Deligne, G.\ Lusztig, \textit{Representations of reductive groups over finite fields}, Announced in Bull.\ Amer.\ Math.\ Soc.\ 80 (1974), 917--921.

\bibitem{SGA4.5} P.\ Deligne, \textit{Cohomologie \'Etale (SGA 4\textonehalf)}, Lecture Notes in Math.\ 569, Springer-Verlag, 1977.

\bibitem{SGA4} M.\ Artin et al., \textit{Th\'eorie des topos et cohomologie \'etale des sch\'emas (SGA 4)}, Lecture Notes in Math.\ 269, 270, 305, Springer-Verlag, 1972--1973.

\bibitem{Green} J.\ A.\ Green, \textit{The characters of the finite general linear groups}, Trans.\ Amer.\ Math.\ Soc.\ 80 (1955), 402--447.

\bibitem{Gro} A.\ Grothendieck, \textit{Formule de Lefschetz et rationalit\'e des fonctions $L$}, S\'eminaire Bourbaki 279, 1964--65.

\bibitem{Han} H.\ C.\ Hansen, \textit{On cycles in flag manifolds}, Math.\ Scand.\ 33 (1973), 269--274.

\bibitem{Mac} I.\ G.\ Macdonald, \textit{The Hall algebra of a symmetric function}, in preparation.

\bibitem{SK} T.\ A.\ Springer, R.\ Steinberg, \textit{Conjugacy classes}, Seminar on algebraic groups and related finite groups, Lecture Notes in Math.\ 131, Springer, 1970, 167--266.

\bibitem{Sri} B.\ Srinivasan, \textit{The characters of the finite symplectic group $Sp(4,q)$}, Trans.\ Amer.\ Math.\ Soc.\ 131 (1968), 488--525.

\bibitem{Stein} R.\ Steinberg, \textit{Endomorphisms of linear algebraic groups}, Mem.\ Amer.\ Math.\ Soc.\ 80 (1968).

\end{thebibliography}

\end{document}
